% concatenated from main_22_.tex
\documentclass[11pt]{amsart}

\usepackage{amsmath,amssymb,amsfonts,amsthm}
\usepackage{mathtools}
\usepackage{amsmath}
\usepackage{graphics}
\usepackage{tikz-cd}
%\usepackage{refcheck}
%\usepackage{cjm2021doi}
%\usepackage[backref]{hyperref}
\newtheorem{theorem}{Theorem}[section]
\newtheorem{corollary}[theorem]{Corollary}
\newtheorem{lemma}[theorem]{Lemma}
\newtheorem{proposition}[theorem]{Proposition}

\newtheorem{definition}[theorem]{Definition}
\newtheorem{remark}[theorem]{Remark}
\newtheorem{example}[theorem]{Example}
\newcommand{\norm}[1]{\left\Vert#1\right\Vert}
\newcommand{\abs}[1]{\left\vert#1\right\vert}
\newcommand{\set}[1]{\left\{#1\right\}}
\newcommand{\R}{\mathbb R}
\newcommand{\K}{\mathbb{K}}

\newcommand{\N}{\mathbb N}
\newcommand{\Z}{\mathbb Z}
\newcommand{\eps}{\varepsilon}
\newcommand{\To}{\longrightarrow}
\newcommand{\C}{\mathcal{C}}
\newcommand{\fK}{\mathbb{K}}
\newcommand{\F}{\mathcal{F}}
\newcommand{\p}{\mathcal{P}}
\newcommand{\A}{\mathcal{A}}
\newcommand{\D}{\mathcal{D}}
\newcommand{\B}{\mathcal{B}}
\newcommand{\U}{\mathcal{U}}
\newcommand{\rel}{\mathcal{R}}
\newcommand{\id}{Id_{]0;1[}}
\newcommand{\lo}{\mathcal{L}}
\newcommand{\mcc}{\mathcal{C}}
\newcommand{\mch}{\mathcal{H}}
\newcommand{\mbH}{\mathbb{H}}
\title{
On spherical Milnor Classifying Spaces I:\\
differential geometry
}

\author{Jean-Pierre Magnot}

\address{{SFR MATHSTIC, LAREMA, Universit\'e d’Angers, 2 Bd Lavoisier, 
49045 Angers cedex 1, France;  Lyc\'ee Jeanne d'Arc, 40 avenue de Grande Bretagne, 63000 Clermont-Ferrand, 
France}; Lepage Research Institute, 17 novembra 1, 081 16 Presov, Slovakia}
\email{\small magnot@math.cnrs.fr; jean-pierr.magnot@ac-clermont.fr; jp.magnot@gmail.com}

\subjclass[2020]{53C23, 57R19, 58B25, 53C27}
\keywords{Milnor classifying space, diffeology, gerbes, Dirac operator, Riemannian geometry, infinite-dimensional geometry}



\begin{document}

\maketitle

\begin{abstract}
We develop a geometric framework for generalized Milnor classifying spaces in the setting of diffeological spaces and infinite-dimensional geometry.

Starting from Milnor's construction, we introduce spherical and projective models endowed with natural diffeological structures compatible with gluing operations. We then investigate their tangent structures, with particular attention to the behavior at the boundary of simplices and to the distinction between tangential and higher-order normal directions.

A natural Riemannian metric is constructed from a barycentric energy functional, leading to a consistent differential calculus on these spaces. This allows us to define differential forms, a Hodge-type theory based on formal adjoints, and a Laplacian without relying on the Hodge star operator.

We further introduce Clifford structures and Dirac operators adapted to this framework, and study twisted versions associated with $\mathbb{Z}_2$-local systems. These structures naturally relate to non-abelian extensions and gerbe-type geometries in the sense of infinite-dimensional Lie theory.

The results provide a coherent geometric setting combining classifying spaces, diffeology, and higher geometric structures.
\end{abstract}



\section{Introduction}

The construction of classifying spaces plays a central role in geometry and topology, notably in the theory of principal bundles, characteristic classes, and differential forms \cite{Milnor1956I,Milnor1956II,BottTu1982}. Among the various models, Milnor's construction provides a particularly flexible approach, based on infinite joins and barycentric coordinates. Its extension to generalized smooth settings is naturally related to diffeology, introduced by Souriau \cite{Souriau1980} and systematically developed by Iglesias-Zemmour \cite{IglesiasZemmour2013}, as well as to other categories of smooth spaces and convenient frameworks \cite{BaezHoffnungRogers2010,Stacey2011}.

In the diffeological setting, classifying spaces and their smooth structures have been studied from several complementary viewpoints. The \(D\)-topology and smooth classifying spaces were investigated in \cite{ChristensenSinnamonWu2014,ChristensenWuSmooth}, while tangent constructions and tangent bundles for diffeological spaces were developed in \cite{ChristensenWu2016}. The diffeology of Milnor's classifying space was analyzed in \cite{MagnotWatts2017}, and diffeological covariant derivatives were studied in \cite{MagnotConnections}. Related aspects of diffeological geometry, including global and differential structures, also appear in \cite{GMW2023,Magnotloops2019,Savelyev2026,Scalbi2024,Watts2012Thesis}.

The present work continues this line of investigation by developing a systematic study of generalized Milnor classifying spaces endowed with additional geometric structures. The key idea is to replace the usual simplicial barycentric coordinates by spherical or projective coordinates, so that the normalization becomes quadratic. This point of view is also naturally related to information-geometric constructions, where square-root coordinates and Fisher--Rao type metrics play a central role \cite{AmariNagaoka2000,AyJostLeSchwachhofer2017,Rao1945}. The relationship of these constructions with information geometry is outlines in an appendix.

A central part of the paper is devoted to differential-geometric structures on these generalized Milnor spaces. We construct a natural Riemannian metric derived from a barycentric energy functional. We then develop a Hodge-type formalism using adjoints rather than relying primarily on the Hodge star operator.
We also introduce Clifford structures and Dirac operators adapted to the effective tangent pseudo-bundles of the spherical Milnor spaces. These constructions are meant to provide a first-order differential framework compatible with the diffeological and stratified nature of the spaces considered here. 

Another important aspect concerns twisted and higher geometric structures. The spherical and projective constructions naturally lead to \(\mathbb Z_2\)-actions, local systems, and twisted objects. Such structures are closely related to extensions of infinite-dimensional Lie groups \cite{Neeb2004,Neeb2007}, bundle gerbes and non-abelian gerbes \cite{Breen2008,BreenMessing2005,Murray1996,NikolausWaldorf2013,Waldorf2007}, and higher categorical or higher bundle structures \cite{BaezLauda2004,Wockel2008}. They also connect with index-theoretic and field-theoretic constructions involving gerbes and determinant-type objects \cite{CareyMickelssonMurray1997,Freed1995}. Some of these aspects are only developed under their premices here, but will be fully developed in follow-up papers. More precisely, the present paper is the first of a serie of works. It intends to set-up safely the geometrical strucures needed for future developments.
The main results of this article can be summarized as follows:
\begin{itemize}
\item we construct spherical and projective models of Milnor classifying spaces compatible with diffeological structures;
\item we provide a detailed analysis of their tangent structures, including boundary phenomena and higher-order effects;
\item we define a natural Riemannian metric and develop a differential calculus adapted to this setting;
\item we introduce Clifford structures and Dirac operators;
\item we relate the resulting twisted structures to extensions and gerbe-type objects.
\end{itemize}

These results establish a coherent geometric framework combining classifying spaces, diffeology, Riemannian-type structures, Dirac operators, and higher geometric data. They are intended as a foundation for further developments, in particular for the study of integrable systems and geometric flows.

The paper is organized as follows. In Section~2, we recall Milnor classifying spaces and their diffeological counterparts. Section~3 is devoted to spherical and projective models. In Section~4, we analyze diffeological smoothness and gluing properties. Section~5 discusses tangent structures and higher-order phenomena. In Section~6, we construct the Riemannian metric and study its non-degeneracy properties. Section~7 develops the differential calculus and Hodge-type theory. Section~8 introduces Clifford structures and Dirac operators. Finally, Section~9 is devoted to twisted structures and their relation with gerbes and section 10 and 11 develops few examples.

	
\section{Diffeological preliminaries and Milnor classifying spaces}
\label{sec:diff-prelim}

\subsection{Diffeological spaces}

We recall the basic notions from diffeology that will be used throughout the paper. General references are \cite{IglesiasZemmour2013} and the survey \cite{GMW2023}. Tangent structures are mainly taken from \cite{ChristensenWu2016,GMW2023}, while pseudo-bundles follow \cite{Pervova2016,Magnot-Cauchy,MagnotConnections}. Related references include \cite{Pervova2017,Pervova2019}.

Diffeology provides a flexible framework for generalized smooth spaces by replacing atlases with families of parametrizations. This makes it possible to treat manifolds, infinite-dimensional spaces, quotients, and singular spaces in a common smooth category.

\begin{definition}[Diffeology]
Let \(X\) be a set. A \emph{parametrization} of \(X\) is a map
\[
p:U\to X,
\]
where \(U\) is an open subset of some Euclidean space. A \emph{diffeology} \(\mathcal P\) on \(X\) is a set of parametrizations, called \emph{plots}, satisfying:
\begin{enumerate}
\item constant maps are plots;
\item plots satisfy the locality condition;
\item plots are stable under smooth reparametrization.
\end{enumerate}
A set equipped with a diffeology is called a \emph{diffeological space}.
\end{definition}

\begin{definition}[Smooth maps]
Let \((X,\mathcal P_X)\) and \((Y,\mathcal P_Y)\) be diffeological spaces. A map
\[
F:X\to Y
\]
is \emph{smooth} if \(F\circ p\in\mathcal P_Y\) for every plot \(p\in\mathcal P_X\).
\end{definition}

The category of diffeological spaces is complete, cocomplete, and forms a quasi-topos \cite{BaezHoffnungRogers2010}. We shall use standard constructions such as pull-back, push-forward, product, subset, and quotient diffeologies \cite{Sou,IglesiasZemmour2013}.

\begin{definition}[Pull-back and push-forward diffeologies]
Let \(f:X\to X'\) be a map.
\begin{itemize}
\item If \(X'\) is diffeological, the pull-back diffeology on \(X\) is
\[
f^*(\mathcal P')=\{p:O_p\to X\mid f\circ p\in\mathcal P'\}.
\]
\item If \(X\) is diffeological, the push-forward diffeology on \(X'\) is the coarsest diffeology making all maps \(f\circ p\) plots.
\end{itemize}
A map \(f:X\to X'\) is called a \emph{subduction} if the diffeology of \(X'\) is the push-forward of the diffeology of \(X\).
\end{definition}

\begin{definition}[Product, subset and quotient diffeologies]
Products are endowed with the coarsest diffeology making all projections smooth. Subsets are endowed with the subset diffeology. Quotients are endowed with the push-forward diffeology induced by the quotient projection.
\end{definition}

When \(X\) is a finite- or infinite-dimensional smooth manifold modeled on a complete locally convex space, it carries its natural \emph{nebulae diffeology}, consisting of the usual smooth parametrizations from open subsets of Euclidean spaces.

\subsection{Frölicher spaces}

We shall occasionally use Frölicher spaces as an intermediate framework between manifolds and diffeological spaces; see \cite{KM,Ma2006-3,Watts2012Thesis,BIgKWa2014}.

\begin{definition}[Frölicher space]
A \emph{Frölicher space} is a triple
\[
(X,\mathcal F,\mathcal C),
\]
where \(\mathcal C\) is a set of contours \(c:\mathbb R\to X\) and \(\mathcal F\) is a set of functions \(f:X\to\mathbb R\), determined by the mutual conditions:
\[
f\in\mathcal F \iff f\circ c\in C^\infty(\mathbb R,\mathbb R)
\quad\text{for all } c\in\mathcal C,
\]
and
\[
c\in\mathcal C \iff f\circ c\in C^\infty(\mathbb R,\mathbb R)
\quad\text{for all } f\in\mathcal F.
\]
\end{definition}

Every Frölicher space carries a natural diffeology \(\mathcal P_\infty(\mathcal F)\), whose plots are the maps \(p:O\to X\) such that \(f\circ p\) is smooth for all \(f\in\mathcal F\).

\begin{proposition}[\cite{Ma2006-3}]
A map between Frölicher spaces is smooth in the Frölicher sense if and only if it is smooth for the associated nebulae diffeologies.
\end{proposition}

Thus one has the informal hierarchy
\[
\text{smooth manifolds}
\;\Longrightarrow\;
\text{Frölicher spaces}
\;\Longrightarrow\;
\text{diffeological spaces}.
\]

Conversely, every diffeological space admits a canonical Frölicher completion generated by its smooth real-valued functions \cite{Watts2012Thesis,BIgKWa2014}. In a Fréchet setting, Frölicher smoothness, \(c^\infty\)-smoothness, and Gâteaux smoothness coincide \cite[Theorem 4.11]{KM}.

\subsection{Diffeological groups and pseudo-bundles}

A \emph{diffeological vector space} is a vector space equipped with a diffeology for which addition and scalar multiplication are smooth. A \emph{diffeological group} is a group equipped with a diffeology for which multiplication and inversion are smooth \cite{IglesiasZemmour2013,Pervova2016}.

\begin{definition}[Diffeological fiber pseudo-bundle]
Let \(E\) and \(X\) be diffeological spaces. A smooth surjective map
\[
\pi:E\to X
\]
is called a \emph{diffeological fiber pseudo-bundle} if \(\pi\) is a subduction.
\end{definition}

This notion generalizes both diffeological vector pseudo-bundles \cite{Pervova2016} and Souriau's quantum structures \cite{Sou}. No local triviality is assumed, and no fixed typical fiber is required.

\begin{definition}[Vector pseudo-bundle]
Let \(\mathbb K\) be a diffeological field. A diffeological fiber pseudo-bundle
\[
\pi:E\to X
\]
is a \(\mathbb K\)-vector pseudo-bundle if each fiber \(E_x\) is a diffeological \(\mathbb K\)-vector space and if fiberwise addition and scalar multiplication are smooth.
\end{definition}

When all fibers are isomorphic as diffeological vector spaces, we shall call \(E\to X\) a diffeological vector bundle.

\subsection{Tangent and cotangent structures}

There are several non-equivalent notions of tangent space in diffeology. The review \cite{GMW2023} identifies five main constructions, originating either from germs of paths or from derivations of smooth functions.

The \emph{internal tangent cone} is defined using germs of paths and generalizes the kinematic tangent construction \cite{DN2007-1,KM}. Its vector-space completion gives the \emph{internal tangent space}, studied for diffeological spaces in \cite{He1995,ChristensenSinnamonWu2014}. Another construction, the \emph{external tangent space}, is defined using derivations \cite{Ma2013,GW2020}. These tangent spaces coincide for finite-dimensional manifolds but differ in general diffeological spaces \cite{ChristensenWu2016,GMW2023}.

The cotangent space and the algebra of differential forms are less ambiguous: differential forms are defined plotwise, by assigning to each plot a classical differential form on its domain, compatibly with smooth reparametrizations \cite{IglesiasZemmour2013,Ma2013}.

\subsection{Diffeological Lie groups}

For a diffeological group \(G\), the internal tangent cone at the identity is often a vector space in the classical examples. However, the existence of a canonical Lie algebra structure in full generality remains subtle. We shall work with diffeological groups for which the tangent space at the identity is equipped with a compatible Lie algebra structure; these will be called \emph{diffeological Lie groups}. For technical aspects, see \cite{Les}.

\subsection{Differential forms with values in pseudo-bundles}

Let \(\pi:E\to X\) be a diffeological vector pseudo-bundle. An \(E\)-valued differential \(n\)-form on \(X\) is defined plotwise: for every plot \(p:O_p\to X\), one assigns a smooth fiberwise alternating map
\[
\alpha_p:\wedge^n TO_p\to p^*E,
\]
compatible with smooth reparametrizations and with changes of plots. The resulting space is denoted
\[
\Omega^n(X,E),
\qquad
\Omega(X,E)=\bigoplus_{n\geq 0}\Omega^n(X,E).
\]

The space \(\Omega^n(X,E)\) is endowed with the functional diffeology: a parametrized family of forms is smooth if its evaluation on every plot is smooth. This generalizes the usual diffeology on scalar differential forms \cite{Ma2013}.

One must be careful with wedge products and exterior differentials in this setting. A wedge product requires a smooth fiberwise multiplication on \(E\), and a natural exterior differential for \(E\)-valued forms generally requires additional connection data.

\subsection{Connections and covariant derivatives}

For principal bundles, a connection is classically described by a connection form satisfying the usual equivariance condition. In the diffeological setting, such connection forms can be defined on suitable frame bundles, but this requires care because vector pseudo-bundles need not have a fixed typical fiber and because the diffeology of the frame bundle may depend on the chosen construction.

For this reason, we shall mainly use covariant derivatives, defined plotwise.

\begin{definition}[Covariant derivative]
Let \(\pi:E\to X\) be a diffeological vector pseudo-bundle. A \emph{covariant derivative} on \(E\) is a family
\[
(p^*\nabla)_{p\in\mathcal P_X}
\]
such that for every plot \(p:O_p\to X\), \(p^*\nabla\) is a covariant derivative on the pullback bundle \(p^*E\to O_p\), and these covariant derivatives are compatible with smooth reparametrizations of plots.
\end{definition}

We denote by \(\mathcal{CD}(E)\) the space of covariant derivatives on \(E\).

\begin{theorem}
If \(\mathcal{CD}(E)\) is non-empty, then it is a diffeological affine space modeled on
\[
\Omega^1(X,\operatorname{End}(E)).
\]
\end{theorem}

The bundle \(\operatorname{End}(E)\) is defined fiberwise by smooth linear endomorphisms of the fibers, with a diffeology ensuring smooth evaluation. This construction follows the approach developed in \cite{Magnot-Cauchy,MagnotConnections}.

\subsection{Pseudo-metrics and Riemannian structures}

Since tangent objects in diffeology are often pseudo-bundles rather than genuine vector bundles, Riemannian structures are naturally formulated as pseudo-metrics.

\begin{definition}
A \emph{pseudo-metric} on a diffeological vector pseudo-bundle \(V\to X\) is a smooth symmetric fiberwise bilinear form
\[
g:V\times_X V\to\mathbb R.
\]
\end{definition}

This is the standard substitute for Riemannian metrics in diffeology and is used in the study of Levi-Civita connections and de Rham-type operators \cite{Pervova2017,GW2020}.

Given a pseudo-metric \(g\) and a covariant derivative \(\nabla\), one defines covariant differentiation of sections along tangent directions whenever the chosen tangent model permits such a contraction. In the present article, the relevant metrics and connections will first be constructed on finite-dimensional local models and then transported through the diffeological gluing procedure. This avoids ambiguity at the level of abstract tangent functors.

\subsection{Milnor classifying spaces in diffeology}

We now recall the diffeological Milnor construction and its main properties.

\subsubsection{The Milnor total space}

Let \(G\) be a diffeological group. Define
\[
E_{\mathrm{Mil}}(G):=
\left\{(t_i,g_i)_{i\in\mathbb N}\ \middle|\
t_i\ge 0,\ \sum_i t_i=1,\ g_i\in G,\ \text{and }t_i=0\text{ for }i\gg 0
\right\}.
\]

For each finite subset \(I\subset \mathbb N\), let
\[
\Delta_I=\left\{(t_i)_{i\in I}\in \mathbb R^I\ \middle|\ t_i\ge 0,\ \sum_{i\in I}t_i=1\right\},
\]
endowed with the subset diffeology, and define
\[
\phi_I:\Delta_I\times G^I\to E_{\mathrm{Mil}}(G)
\]
by extending the coordinates by zero outside \(I\).

\begin{definition}
The \emph{Milnor diffeology} on \(E_{\mathrm{Mil}}(G)\) is the final diffeology generated by the family \((\phi_I)_I\).
\end{definition}

This is precisely the diffeology used in \cite{MagnotWatts2017}. Note that the boundary behavior is encoded by the face inclusions \(\Delta_J\hookrightarrow \Delta_I\) for \(J\subset I\): when some \(t_i\) vanish, one passes smoothly to lower-dimensional simplices. :contentReference[oaicite:8]{index=8}

\subsubsection{The classifying space and the universal bundle}

The group \(G\) acts on \(E_{\mathrm{Mil}}(G)\) by
\[
h\cdot (t_i,g_i)_{i\in\mathbb N}=(t_i,hg_i)_{i\in\mathbb N}.
\]

\begin{definition}
The \emph{Milnor classifying space} is the quotient
\[
B_{\mathrm{Mil}}(G):=E_{\mathrm{Mil}}(G)/G
\]
equipped with the quotient diffeology.
\end{definition}

The main results of Magnot--Watts can be summarized as follows.

\begin{theorem}[Magnot--Watts]\label{thm:MW}
Let \(G\) be a diffeological group.
\begin{enumerate}
\item The space \(E_{\mathrm{Mil}}(G)\) is diffeologically contractible.
\item The projection
\[
E_{\mathrm{Mil}}(G)\to B_{\mathrm{Mil}}(G)
\]
is a principal diffeological \(G\)-bundle.
\item For a diffeological space \(X\), smooth homotopy classes of smooth maps
\[
X\to B_{\mathrm{Mil}}(G)
\]
classify \(D\)-numerable principal diffeological \(G\)-bundles over \(X\).
\end{enumerate}
\end{theorem}

This is the precise analogue, in the diffeological category, of the classical Milnor classification theorem. :contentReference[oaicite:9]{index=9}

\subsubsection{Canonical connection}

Another crucial point for the present article is that the Milnor bundle carries a natural connection. In Magnot--Watts, this connection is written by barycentric averaging of the Maurer--Cartan forms of the factors.

\begin{proposition}[Magnot--Watts]
The principal bundle
\[
E_{\mathrm{Mil}}(G)\to B_{\mathrm{Mil}}(G)
\]
carries a canonical connection whose local connection form is given by a weighted sum
\[
\theta=\sum_i t_i\,\theta_i,
\]
where \(\theta_i\) denotes the Maurer--Cartan form pulled back from the \(i\)-th copy of \(G\).
\end{proposition}

This canonical connection is one of the main reasons why the Milnor model is especially well suited for the geometric constructions developed later in the paper. :contentReference[oaicite:10]{index=10}

\subsubsection{Comparison with Christensen--Wu}

Christensen and Wu develop a parallel theory of smooth classifying spaces for diffeological groups and prove that \(D\)-numerable smooth principal bundles are classified by smooth maps into their classifying space. Their construction is not based on the explicit Milnor join model, but on a general smooth homotopy-theoretic approach to \(D\)-numerable bundles. :contentReference[oaicite:11]{index=11}

\begin{remark}
For the purposes of the present article, the two approaches play different roles:
\begin{itemize}
\item Christensen--Wu provide a very general classification framework for smooth principal bundles;
\item the Magnot--Watts Milnor model provides explicit barycentric coordinates and a canonical connection.
\end{itemize}
It is this second feature that makes the Milnor model particularly suitable for the introduction of spherical/projective coordinates, metrics, Clifford structures and Dirac operators.
\end{remark}

\subsubsection{What is retained for the present work}

From the preceding results, we retain the following facts:
\begin{enumerate}
\item the Milnor total space is built by final diffeology from finite-dimensional simplex charts;
\item its boundary behavior is controlled by face inclusions and thus by vanishing barycentric coordinates;
\item it carries a universal principal bundle structure and a canonical barycentric connection;
\item it admits an explicit quotient model suitable for geometric refinements.
\end{enumerate}

These are exactly the ingredients that will be modified in the next section, where we replace the simplicial normalization \(\sum_i t_i=1\) by a spherical normalization \(\sum_i x_i^2=1\).
\section{Spherical and projective models} \label{3}

\subsection{From simplicial to spherical coordinates}

Recall that in Milnor's construction the barycentric coordinates satisfy
\[
t_i \ge 0, \qquad \sum_i t_i = 1,
\]
and that the diffeology on the simplices $\Delta_I$ is not the standard subset diffeology, but is adapted so that plots are stationary at the boundary, see~\cite{MagnotWatts2017}.

In particular, if a plot $u \mapsto (t_i(u))$ reaches a boundary point where some $t_i(u)=0$, then the corresponding coordinate is locally constant.

\medskip

We now introduce spherical coordinates by setting
\[
t_i = x_i^2,
\qquad \sum_i x_i^2 = 1.
\]

However, in contrast with the simplicial case, the spheres
\[
S_I = \left\{ (x_i)_{i\in I} \in \mathbb{R}^I \;\middle|\; \sum_{i\in I} x_i^2 = 1 \right\}
\]
are endowed with the subset diffeology induced from $\mathbb{R}^I$.

\begin{remark}
In particular, no stationarity condition is imposed when a coordinate $x_i$ vanishes. This creates a fundamental asymmetry between the simplicial and spherical models.
\end{remark}

\begin{remark}
As a consequence, the map
\[
t_i \mapsto \sqrt{t_i}
\]
is smooth for the Milnor diffeology, while the inverse map
\[
x_i \mapsto x_i^2
\]
is not smooth in general.
\end{remark}

\subsection{Spherical Milnor space and group actions}

\begin{definition}
Define
\[
\widetilde E^{\mathrm{sph}}(G)
:=
\left\{
(x_i,g_i)_{i\in\mathbb{N}}
\;\middle|\;
\sum_i x_i^2 = 1,\ 
x_i \in \mathbb{R},\ 
g_i \in G,\ 
\text{finite support}
\right\}
\]
and introduce the equivalence relation
\[
(x_i,g_i)\sim (x_i,g_i')
\quad\Longleftrightarrow\quad
g_i=g_i' \ \text{whenever } x_i\neq 0.
\]
\end{definition}

\begin{definition}
The \emph{spherical Milnor space} is the quotient
\[
E^{\mathrm{sph}}(G) := \widetilde E^{\mathrm{sph}}(G)/\sim.
\]
\end{definition}

\begin{remark}
This relation ensures that group coordinates corresponding to vanishing weights are irrelevant, in exact analogy with Milnor's construction.
\end{remark}

\medskip

For each finite $I \subset \mathbb{N}$, define
\[
S_I := \left\{ (x_i)_{i\in I} \in \mathbb{R}^I \;\middle|\; \sum_{i\in I} x_i^2 = 1 \right\}.
\]

\begin{definition}
The diffeology on $E^{\mathrm{sph}}(G)$ is the final diffeology generated by the maps
\[
\psi_I : S_I \times G^I \to E^{\mathrm{sph}}(G),
\]
defined by extension by zero and passage to the quotient.
\end{definition}

\begin{remark}
This diffeology is compatible with the equivalence relation, since only coordinates with $x_i \neq 0$ contribute.
\end{remark}

\subsubsection{Group actions}

\begin{definition}
The group $G$ acts on $E^{\mathrm{sph}}(G)$ by
\[
h \cdot (x_i,g_i) = (x_i, hg_i).
\]
\end{definition}

\begin{proposition}
This action is well-defined and smooth.
\end{proposition}

\begin{proof}
If $x_i=0$, the coordinate $g_i$ is irrelevant, hence the action is compatible with the equivalence relation. Smoothness follows from the local model $S_I \times G^I$.
\end{proof}

\begin{definition}
The group $\mathbb{Z}_2$ acts on $E^{\mathrm{sph}}(G)$ by
\[
\varepsilon \cdot (x_i,g_i) = (\varepsilon x_i,g_i),
\qquad \varepsilon \in \{\pm 1\}.
\]
\end{definition}

\begin{proposition}
This action is well-defined and smooth.
\end{proposition}

\begin{proof}
It acts only on the spherical coordinates and preserves the constraint $\sum x_i^2=1$. Smoothness is clear on each $S_I$.
\end{proof}

\begin{remark}
More generally, one may consider independent sign changes $(\mathbb{Z}_2)^I$, but in the present work we restrict to the global $\mathbb{Z}_2$-action.
\end{remark}

\subsubsection{Quotient spaces}

\begin{definition}
We define the following quotient spaces:
\[
E^{\mathrm{sph}}(G)/G,
\qquad
E^{\mathrm{sph}}(G)/\mathbb{Z}_2,
\qquad
B^{\mathrm{sph}}(G) := E^{\mathrm{sph}}(G)/(G \times \mathbb{Z}_2).
\]
\end{definition}

\begin{proposition}
The actions of $G$ and $\mathbb{Z}_2$ commute.
\end{proposition}

\begin{proof}
We have
\[
h \cdot (\varepsilon \cdot (x_i,g_i))
= h \cdot (\varepsilon x_i, g_i)
= (\varepsilon x_i, hg_i),
\]
and
\[
\varepsilon \cdot (h \cdot (x_i,g_i))
= \varepsilon \cdot (x_i, hg_i)
= (\varepsilon x_i, hg_i),
\]
hence the actions commute.
\end{proof}

\begin{corollary}
There are canonical identifications
\[
(E^{\mathrm{sph}}(G)/G)/\mathbb{Z}_2
\simeq
(E^{\mathrm{sph}}(G)/\mathbb{Z}_2)/G
\simeq
B^{\mathrm{sph}}(G).
\]
\end{corollary}

\begin{remark}
Thus the quotients by $G$ and $\mathbb{Z}_2$ can be performed in any order.
\end{remark}

\subsubsection{Interpretation}

\begin{remark}
The space $E^{\mathrm{sph}}(G)$ plays the role of a universal space endowed with an additional $\mathbb{Z}_2$-symmetry.

The quotient $E^{\mathrm{sph}}(G)/G$ corresponds to a spherical version of the classifying space, while $E^{\mathrm{sph}}(G)/\mathbb{Z}_2$ corresponds to a projectivized version.

The double quotient $B^{\mathrm{sph}}(G)$ provides a refined classifying space incorporating both symmetries.
\end{remark}
\subsection{Embedding of the Milnor space}

\begin{definition}
Define
\[
\iota : E(G) \longrightarrow E^{\mathrm{sph}}(G),
\qquad
(t_i,g_i) \longmapsto (\sqrt{t_i},g_i).
\]
\end{definition}

\begin{proposition}
The map $\iota$ is well-defined.
\end{proposition}

\begin{proof}
If $t_i = 0$, the coordinate $g_i$ is irrelevant in $E(G)$, and similarly in $E^{\mathrm{sph}}(G)$ the corresponding coordinate $x_i=0$ makes $g_i$ irrelevant. Hence the definition is consistent with both quotient structures.
\end{proof}

\begin{proposition}
The map $\iota$ is smooth.
\end{proposition}

\begin{proof}
Let $p:U\to E(G)$ be a plot. By definition of the Milnor diffeology, locally
\[
p(u) = (t_i(u),g_i(u))
\]
with $(t_i(u))$ factoring through a simplex $\Delta_I$ and such that each $t_i$ is smooth and stationary at the boundary.

Then
\[
\iota \circ p(u) = (\sqrt{t_i(u)},g_i(u)).
\]

Since $t_i$ is stationary at $0$, the function $\sqrt{t_i(u)}$ is smooth. Hence the map
\[
u \mapsto (\sqrt{t_i(u)},g_i(u))
\]
is smooth as a map into $S_I \times G^I$, and therefore defines a plot of $E^{\mathrm{sph}}(G)$.
\end{proof}

\subsection{Non-smoothness of the inverse map}

\begin{proposition}
The map
\[
\Phi : E^{\mathrm{sph}}(G) \to E(G),
\qquad
(x_i,g_i) \mapsto (x_i^2,g_i)
\]
is not smooth in general.
\end{proposition}

\begin{proof}
Consider a plot $(x_i(u))$ in $E^{\mathrm{sph}}(G)$ such that some coordinate $x_i(u)$ crosses $0$ with nonzero derivative. Then $x_i(u)^2$ is not stationary at $0$, hence does not define a plot in the Milnor diffeology on $E(G)$.
\end{proof}

\begin{remark}
This reflects the asymmetry of the Milnor--Watts diffeology: the map $t \mapsto \sqrt{t}$ is smooth, while $x \mapsto x^2$ is not.
\end{remark}

\subsection{Contractibility}

\begin{theorem}
The space $E^{\mathrm{sph}}(G)$ is diffeologically contractible.
\end{theorem}

\begin{proof}
We construct a contraction using the shift operator.

Define
\[
S(x_1,x_2,x_3,\dots) = (0,x_1,x_2,\dots),
\]
and extend it to $E^{\mathrm{sph}}(G)$ by inserting the identity element of $G$ at the first coordinate.

Define
\[
H_t^{(1)}(x) = \frac{(1-t)x + t Sx}{\|(1-t)x + t Sx\|}.
\]
Since $x$ and $Sx$ have disjoint supports, they are orthogonal, hence the denominator never vanishes.

This defines a smooth homotopy from the identity to the shift.

Then define
\[
H_t^{(2)}(y) = \frac{(1-t)Sy + t e_1}{\|(1-t)Sy + t e_1\|},
\]
where $e_1 = (1,0,0,\dots)$.

This defines a contraction of the shift to the base point.

Composing the two homotopies yields a contraction of $E^{\mathrm{sph}}(G)$.
\end{proof}

\subsection{Geometric interpretation}

\begin{remark}
The spherical normalization replaces the affine constraint
\[
\sum_i t_i = 1
\]
by a quadratic constraint
\[
\sum_i x_i^2 = 1,
\]
which is invariant under orthogonal transformations.
\end{remark}

\begin{remark}
The embedding $\iota$ identifies $E(G)$ with a distinguished subspace of $E^{\mathrm{sph}}(G)$ adapted to the Milnor diffeology, while the ambient spherical space is better suited for metric and Dirac-type constructions.
\end{remark}

\subsection{Summary}

The spherical Milnor space:
\begin{itemize}
\item is a quotient space compatible with vanishing coordinates,
\item carries a natural diffeological structure,
\item contains $E(G)$ as a smooth subspace,
\item and is diffeologically contractible.
\end{itemize}
\section{Diffeological smoothness and gluing} \label{4}

\subsection{Final diffeology and local models}

Recall that $E^{\mathrm{sph}}(G)$ is endowed with the final diffeology generated by the maps
\[
\psi_I : S_I \times G^I \longrightarrow E^{\mathrm{sph}}(G),
\]
where $I \subset \mathbb{N}$ is finite.

\begin{definition}
A map $p : U \to E^{\mathrm{sph}}(G)$ is a plot if and only if for every $u \in U$ there exist:
\begin{itemize}
\item an open neighborhood $V \subset U$ of $u$,
\item a finite index set $I$,
\item a smooth map $\widetilde p : V \to S_I \times G^I$,
\end{itemize}
such that
\[
p|_V = \psi_I \circ \widetilde p.
\]
\end{definition}

\subsection{Directed system and gluing maps}

For $J \subset I$, define the inclusion
\[
\iota_{J,I} : S_J \times G^J \longrightarrow S_I \times G^I
\]
by extending coordinates by zero and identities.

\begin{proposition}
The family $(S_I \times G^I, \iota_{J,I})$ forms a directed system in the category of diffeological spaces.
\end{proposition}

\begin{proposition}
The space $E^{\mathrm{sph}}(G)$ is the colimit of this system in the category of diffeological spaces.
\end{proposition}

\begin{proof}
By construction, the diffeology on $E^{\mathrm{sph}}(G)$ is the final diffeology generated by the maps $\psi_I$, which precisely encodes the colimit structure.
\end{proof}

\subsection{Local finiteness}

\begin{proposition}
Every plot $p : U \to E^{\mathrm{sph}}(G)$ is locally supported on a fixed finite index set $I$.
\end{proposition}

\begin{proof}
Immediate from the definition of the final diffeology.
\end{proof}

\subsection{Compatibility under change of charts}

\begin{proposition}
Let $p : U \to E^{\mathrm{sph}}(G)$ be a plot. If on overlapping neighborhoods $V_1, V_2 \subset U$ the map $p$ factors through
\[
S_{I_1} \times G^{I_1}, \qquad S_{I_2} \times G^{I_2},
\]
then on $V_1 \cap V_2$ it factors through
\[
S_{I_1 \cup I_2} \times G^{I_1 \cup I_2}.
\]
\end{proposition}

\begin{proof}
This follows from the compatibility of the inclusions $\iota_{J,I}$.
\end{proof}

\subsection{Stratification by support}

\begin{definition}
For each finite $I \subset \mathbb{N}$, define
\[
E_I := \psi_I(S_I \times G^I).
\]
\end{definition}

\begin{proposition}
Each $E_I$ is a diffeological subspace of $E^{\mathrm{sph}}(G)$ and is diffeomorphic to $S_I \times G^I$.
\end{proposition}

\begin{proposition}
We have
\[
E^{\mathrm{sph}}(G) = \bigcup_{I} E_I.
\]
\end{proposition}

\begin{remark}
The family $(E_I)$ defines a stratification indexed by finite subsets of $\mathbb{N}$.
\end{remark}

\subsection{Behavior at the boundary}

\begin{proposition}
Let $x \in E_I$ such that some coordinate $x_j = 0$. Then $x$ lies in the image of $E_J$ for some strict subset $J \subset I$.
\end{proposition}

\begin{proof}
If $x_j = 0$, the support of $(x_i)$ is contained in a strict subset $J \subset I$.
\end{proof}

\begin{remark}
Thus, the boundary of each stratum is contained in lower-dimensional strata.
\end{remark}

\subsection{Smooth functions}

\begin{proposition}
A function $f : E^{\mathrm{sph}}(G) \to \mathbb{R}$ is smooth if and only if for every finite $I$, the pullback
\[
f \circ \psi_I : S_I \times G^I \to \mathbb{R}
\]
is smooth.
\end{proposition}

\begin{proof}
This follows directly from the definition of the final diffeology.
\end{proof}

\subsection{Summary}

The diffeological structure on $E^{\mathrm{sph}}(G)$ is characterized by:
\begin{itemize}
\item a colimit structure over finite-dimensional smooth spaces,
\item local finite-dimensionality of plots,
\item compatibility under inclusions of index sets,
\item and a stratified structure governed by supports of coordinates.
\end{itemize}
\section{Tangent structures and boundary behavior} \label{5}

\subsection{Local tangent model}

Let $E^{\mathrm{sph}}(G)$ be as in Section~3.

\begin{proposition}
Locally, near any point of support $I$, the space $E^{\mathrm{sph}}(G)$ is modeled on
\[
S_I \times G^I.
\]
\end{proposition}

\begin{definition}
The tangent space at a point $(x_i,g_i)$ with support $I$ is defined locally as
\[
T_{(x,g)} E^{\mathrm{sph}}(G)
\simeq T_x S_I \oplus \bigoplus_{i\in I} T_{g_i} G.
\]
\end{definition}

\subsection{Tangent space to the sphere}

Recall that
\[
T_x S_I = \left\{ v \in \mathbb{R}^I \;\middle|\; \sum_{i\in I} x_i v_i = 0 \right\}.
\]

\begin{remark}
This is a codimension-one subspace of $\mathbb{R}^I$.
\end{remark}

\subsection{Behavior at the boundary}

Let $(x_i,g_i)$ be such that some coordinate $x_j = 0$.

\begin{proposition}
At such a point, the tangent space satisfies
\[
v_j \ \text{is unconstrained by the condition } \sum x_i v_i = 0.
\]
\end{proposition}

\begin{proof}
The constraint reads
\[
\sum_i x_i v_i = 0.
\]
If $x_j = 0$, then $v_j$ does not appear in the equation.
\end{proof}

\begin{remark}
This shows that directions corresponding to vanishing coordinates are tangent directions.
\end{remark}

\subsection{Absence of normal directions at first order}

\begin{proposition}
There is no distinguished normal direction at first order at boundary points.
\end{proposition}

\begin{proof}
At a point where $x_j=0$, any variation of $x_j$ satisfies the tangent condition. Hence no direction is singled out as transverse.
\end{proof}

\begin{remark}
This contrasts with manifolds with boundary, where a normal direction is canonically defined.
\end{remark}

\subsection{Second-order behavior}

\begin{proposition}
Second-order variations detect normal directions.
\end{proposition}

\begin{proof}
Let $x_j=0$. Consider a path
\[
x_j(t) = t.
\]
Then
\[
t_j(t) = x_j(t)^2 = t^2.
\]

The first derivative vanishes at $t=0$, but the second derivative does not. Hence normal behavior appears at second order.
\end{proof}

\begin{remark}
Thus, normal directions are invisible at the level of the tangent space but appear in higher-order tangent structures.
\end{remark}

\subsection{Interpretation}

\begin{remark}
This phenomenon reflects the quadratic nature of the spherical constraint and the singular behavior of the map $x \mapsto x^2$ at zero.
\end{remark}

\begin{remark}
It also explains why the spherical model is better suited for metric constructions: it regularizes first-order singularities.
\end{remark}
\subsection{Stratification and first-order invisibility}

We formalize the stratified structure induced by the support of the spherical coordinates and its behavior with respect to tangent structures.

\begin{definition}
For each finite subset $I \subset \mathbb{N}$, define the stratum
\[
E_I := \psi_I(S_I \times G^I) \subset E^{\mathrm{sph}}(G).
\]
\end{definition}

\begin{proposition}
Each $E_I$ is a diffeological subspace of $E^{\mathrm{sph}}(G)$, and the map
\[
\psi_I : S_I \times G^I \longrightarrow E_I
\]
is a diffeomorphism.
\end{proposition}

\begin{proof}
This follows directly from the definition of the quotient and of the final diffeology.
\end{proof}

\begin{proposition}
We have a stratification
\[
E^{\mathrm{sph}}(G) = \bigcup_{I \subset \mathbb{N},\ \text{finite}} E_I,
\]
and for $J \subset I$,
\[
E_J \subset \overline{E_I}.
\]
\end{proposition}

\begin{proof}
The first statement follows from the finite support condition. The second follows from the fact that vanishing coordinates define lower-dimensional strata.
\end{proof}

\begin{theorem}[Stratification is invisible at first order]
Let $x \in E_J$ for some finite $J \subset \mathbb{N}$. Then the tangent space at $x$ satisfies
\[
T_x E^{\mathrm{sph}}(G) \simeq T_x E_I
\]
for any $I \supset J$.

In particular, the inclusion of strata
\[
E_J \hookrightarrow E_I
\]
does not induce any distinguished normal direction at the level of tangent spaces.
\end{theorem}

\begin{proof}
Let $x$ have support $J$. Locally, $x$ is represented in $S_J \times G^J$.

If we embed into a larger chart $S_I \times G^I$, the tangent space is
\[
T_x S_I = \left\{ v \in \mathbb{R}^I \;\middle|\; \sum_{i\in I} x_i v_i = 0 \right\}.
\]

For $i \in I \setminus J$, we have $x_i = 0$, hence no constraint is imposed on $v_i$. Therefore the additional directions corresponding to $I \setminus J$ are tangent directions.

Thus the tangent space does not distinguish between $E_J$ and its ambient $E_I$.
\end{proof}

\begin{corollary}
The stratification of $E^{\mathrm{sph}}(G)$ cannot be detected by first-order differential data.
\end{corollary}

\begin{proposition}[Second-order detection]
Let $x \in E_J$ and consider a curve $x(t)$ such that for some $i \notin J$,
\[
x_i(t) = t.
\]
Then
\\subsection{Support stratification and first-order tangent structure}

We now describe the tangent behavior of the spherical Milnor space along the support strata. The point is subtle: in the spherical model, a vanishing coordinate may be activated at first order, whereas its square is invisible at first order in the simplicial Milnor coordinates.

For a finite subset \(I\subset \mathbb N\), set
\[
Q_I
:=
(S_I\times G^I)/\!\sim_I,
\]
where
\[
(x_i,g_i)_{i\in I}\sim_I (x_i,g_i')_{i\in I}
\quad\Longleftrightarrow\quad
g_i=g_i' \ \text{for every } i\in I \text{ such that } x_i\neq 0.
\]
Thus the group coordinate \(g_i\) is forgotten whenever \(x_i=0\).

The map
\[
S_I\times G^I\longrightarrow E^{\mathrm{sph}}(G)
\]
factors through \(Q_I\), and we denote by
\[
\psi_I:Q_I\longrightarrow E^{\mathrm{sph}}(G)
\]
the induced map.

\begin{definition}
For a point
\[
p=[(x_i,g_i)]\in E^{\mathrm{sph}}(G),
\]
its \emph{support} is
\[
\operatorname{supp}(p):=\{i\in\mathbb N\mid x_i\neq 0\}.
\]
For a finite subset \(J\subset \mathbb N\), define the support stratum
\[
E_J^\circ
:=
\{p\in E^{\mathrm{sph}}(G)\mid \operatorname{supp}(p)=J\}.
\]
\end{definition}

\begin{remark}
The strata \(E_J^\circ\) are indexed by finite subsets \(J\subset\mathbb N\). Their closures meet along strata of smaller support. In particular, when a coordinate \(x_i\) tends to zero, the corresponding group coordinate \(g_i\) becomes irrelevant.
\end{remark}

\begin{lemma}
Let \(p=[(x_i,g_i)]\in E^{\mathrm{sph}}(G)\) have support \(J\), and let \(I\supset J\) be finite. Then the internal tangent space of \(Q_I\) at \(p\) is naturally identified with
\[
T_p Q_I
\simeq
\left\{
(v_i)_{i\in I}\in \mathbb R^I
\;\middle|\;
\sum_{j\in J} x_j v_j=0
\right\}
\oplus
\bigoplus_{j\in J} T_{g_j}G.
\]
No tangent component in \(T_{g_i}G\) occurs for indices \(i\in I\setminus J\).
\end{lemma}

\begin{proof}
Choose a representative of \(p\) in \(S_I\times G^I\). A plot through \(p\) is locally represented by a smooth map
\[
u\longmapsto (x_i(u),g_i(u))_{i\in I}
\]
with
\[
\sum_{i\in I}x_i(u)^2=1.
\]
Differentiating this relation at \(u=0\) gives
\[
2\sum_{i\in I}x_i(0)\dot x_i(0)=0.
\]
Since \(x_i(0)=0\) for \(i\in I\setminus J\), this becomes
\[
\sum_{j\in J}x_j\dot x_j(0)=0.
\]
Thus the spherical part of a tangent vector is an element
\[
(v_i)_{i\in I}
\quad\text{with}\quad
\sum_{j\in J}x_jv_j=0.
\]

For the group part, if \(j\in J\), then \(x_j\neq0\), and the coordinate \(g_j\) is meaningful in the quotient. A curve \(g_j(u)\) therefore contributes a tangent vector in \(T_{g_j}G\).

If \(i\in I\setminus J\), then \(x_i=0\) at the base point, and the quotient relation identifies all possible values of \(g_i\). Hence a curve varying only \(g_i\), with \(x_i\equiv0\), represents the zero tangent direction. More generally, at first order, the variation of \(g_i\) is not independently visible at the point where \(x_i=0\); only the variation of the spherical coordinate \(x_i\) is visible. Therefore no summand \(T_{g_i}G\) appears for \(i\notin J\).

Conversely, given
\[
(v_i)_{i\in I}
\quad\text{with}\quad
\sum_{j\in J}x_jv_j=0
\]
and tangent vectors \(\xi_j\in T_{g_j}G\) for \(j\in J\), choose smooth curves \(g_j(t)\) in \(G\) with
\[
g_j(0)=g_j,\qquad \dot g_j(0)=\xi_j.
\]
Set
\[
y(t)=x+tv\in \mathbb R^I.
\]
Since
\[
\frac{d}{dt}\bigg|_{t=0}\|y(t)\|^2
=
2\sum_{j\in J}x_jv_j=0,
\]
the normalized curve
\[
x(t)=\frac{y(t)}{\|y(t)\|}
\]
is well-defined for small \(t\) and satisfies
\[
x(0)=x,\qquad \dot x(0)=v.
\]
Together with the curves \(g_j(t)\), and arbitrary constant choices of \(g_i\) for \(i\notin J\), this gives a plot realizing the prescribed tangent vector. This proves the identification.
\end{proof}

\begin{theorem}[First-order tangent structure along support strata]
Let \(p=[(x_i,g_i)]\in E^{\mathrm{sph}}(G)\) have support \(J\). Then the internal tangent space of \(E^{\mathrm{sph}}(G)\) at \(p\) is the filtered colimit
\[
T_pE^{\mathrm{sph}}(G)
\simeq
\varinjlim_{I\supset J}
\left[
\left\{
(v_i)_{i\in I}\in \mathbb R^I
\;\middle|\;
\sum_{j\in J}x_jv_j=0
\right\}
\oplus
\bigoplus_{j\in J}T_{g_j}G
\right].
\]
Equivalently,
\[
T_pE^{\mathrm{sph}}(G)
\simeq
\left\{
(v_i)_{i\in\mathbb N}^{\mathrm{fin}}
\;\middle|\;
\sum_{j\in J}x_jv_j=0
\right\}
\oplus
\bigoplus_{j\in J}T_{g_j}G,
\]
where \((v_i)^{\mathrm{fin}}\) denotes a finitely supported real sequence.
\end{theorem}

\begin{proof}
By definition of the final diffeology, every plot of \(E^{\mathrm{sph}}(G)\) through \(p\) locally factors through one of the finite models \(Q_I\) with \(I\supset J\). Therefore the internal tangent space is generated by the tangent spaces \(T_pQ_I\), with compatibility under the inclusions \(I\subset K\).

For \(I\subset K\), the inclusion
\[
Q_I\hookrightarrow Q_K
\]
extends the coordinate vector \((v_i)_{i\in I}\) by zero on \(K\setminus I\), while leaving the group tangent components indexed by \(J\) unchanged. Hence the tangent spaces form a directed system.

Using the previous lemma, the directed colimit is precisely the space of finitely supported sequences \((v_i)\) satisfying
\[
\sum_{j\in J}x_jv_j=0,
\]
together with group tangent components only for active indices \(j\in J\). This gives the stated formula.
\end{proof}

\begin{corollary}
Let \(p\in E_J^\circ\). The tangent space to the stratum \(E_J^\circ\) embeds as
\[
T_pE_J^\circ
\simeq
\left\{
(v_j)_{j\in J}
\;\middle|\;
\sum_{j\in J}x_jv_j=0
\right\}
\oplus
\bigoplus_{j\in J}T_{g_j}G
\]
inside \(T_pE^{\mathrm{sph}}(G)\). The quotient consists of finitely supported variations of the inactive spherical coordinates:
\[
T_pE^{\mathrm{sph}}(G)/T_pE_J^\circ
\simeq
\mathbb R^{(\mathbb N\setminus J)}.
\]
\end{corollary}

\begin{proof}
The tangent space to the stratum is obtained by imposing \(v_i=0\) for all \(i\notin J\). The full tangent space allows finitely supported variations \(v_i\) for all inactive indices \(i\notin J\). Taking the quotient gives the stated direct sum.
\end{proof}

\begin{remark}
This is the precise sense in which the spherical model differs from the simplicial Milnor model. In the spherical model, an inactive coordinate \(x_i=0\) may be activated at first order. However, this first-order activation is symmetric in the sign of \(x_i\); it is not an outward normal direction of a manifold with boundary.
\end{remark}

\begin{theorem}[Quadratic invisibility in Milnor coordinates]
Let \(p\in E_J^\circ\), and let \(v\in T_pE^{\mathrm{sph}}(G)\) be a tangent vector with \(v_i\neq0\) for some inactive index \(i\notin J\). Under the quadratic map
\[
t_i=x_i^2,
\]
the corresponding first-order variation of the Milnor coordinate \(t_i\) vanishes:
\[
\dot t_i(0)=0.
\]
The first nonzero contribution occurs at second order:
\[
\ddot t_i(0)=2v_i^2.
\]
\end{theorem}

\begin{proof}
Let \(x_i(s)\) be a smooth curve with
\[
x_i(0)=0,\qquad \dot x_i(0)=v_i.
\]
Then
\[
t_i(s)=x_i(s)^2.
\]
Differentiating gives
\[
\dot t_i(s)=2x_i(s)\dot x_i(s),
\]
hence
\[
\dot t_i(0)=2x_i(0)\dot x_i(0)=0.
\]
Differentiating again,
\[
\ddot t_i(s)=2\dot x_i(s)^2+2x_i(s)\ddot x_i(s),
\]
and therefore
\[
\ddot t_i(0)=2\dot x_i(0)^2=2v_i^2.
\]
\end{proof}

\begin{remark}
Thus, directions that are visible at first order in the spherical model become second-order phenomena in the simplicial Milnor coordinates. This observation will be important for the comparison between the spherical metric and Fisher--Rao type structures.
\end{remark}]
and the corresponding direction vanishes at first order but not at second order.
\end{proposition}

\begin{proof}
Immediate computation.
\end{proof}

\begin{remark}
Thus, the stratification is only visible at the level of second-order tangent structures.
\end{remark}
\subsection{Summary}

The tangent structure of $E^{\mathrm{sph}}(G)$ exhibits:
\begin{itemize}
\item a standard local description on each stratum,
\item absence of normal directions at first order,
\item emergence of normal behavior at second order.
\end{itemize}
\section{Riemannian structure and energy formulation} \label{6}

\subsection{Local definition of the metric}

Let $G$ be a Lie group endowed with a fixed right-invariant Riemannian metric
\[
\langle \cdot,\cdot \rangle_G.
\]

For a finite index set $I$, consider the local model
\[
S_I \times G^I.
\]

\begin{definition}
On $S_I \times G^I$, define the metric
\[
g_I
=
\sum_{i\in I} dx_i^2
\;+\;
\sum_{i\in I} x_i^2 \,\langle \cdot,\cdot \rangle_G^{(i)},
\]
where $\langle \cdot,\cdot \rangle_G^{(i)}$ denotes the metric on the $i$-th copy of $G$.
\end{definition}

\begin{remark}
The first term is the restriction of the Euclidean metric to the sphere $S_I$, while the second term weights the group directions by the barycentric coefficients $x_i^2$.
\end{remark}

\subsection{Compatibility with the quotient}

\begin{proposition}
The metric $g_I$ is compatible with the equivalence relation defining $E^{\mathrm{sph}}(G)$.
\end{proposition}

\begin{proof}
If $x_i = 0$, then the corresponding group coordinate $g_i$ is irrelevant in the quotient. At such a point, the metric coefficient $x_i^2$ vanishes, hence the contribution of $T_{g_i}G$ to the metric is zero. Therefore the metric descends to the quotient.
\end{proof}

\subsection{Compatibility under inclusions}

\begin{proposition}
For $J \subset I$, the metric $g_I$ restricts to $g_J$ under the natural inclusion
\[
S_J \times G^J \hookrightarrow S_I \times G^I.
\]
\end{proposition}

\begin{proof}
On the inclusion, one has $x_i = 0$ for $i \in I \setminus J$, hence all additional terms vanish. The metric reduces exactly to $g_J$.
\end{proof}

\begin{corollary}
The family $(g_I)_I$ defines a global metric $g$ on $E^{\mathrm{sph}}(G)$.
\end{corollary}

\subsection{Smoothness}

\begin{proposition}
The metric $g$ is smooth in the diffeological sense.
\end{proposition}

\begin{proof}
By definition of diffeological smoothness, it suffices to check that for every plot
\[
p : U \to E^{\mathrm{sph}}(G),
\]
the pullback metric $p^*g$ is smooth.

Locally, $p$ factors through some $S_I \times G^I$, and the metric is smooth on this finite-dimensional manifold. Hence $p^*g$ is smooth.
\end{proof}

\subsection{Behavior at the boundary}

\begin{proposition}
The metric $g$ extends smoothly across strata where some coordinates $x_i$ vanish.
\end{proposition}

\begin{proof}
Near a point where $x_i=0$, the metric term
\[
x_i^2 \,\langle \cdot,\cdot \rangle_G
\]
vanishes smoothly. Therefore no singularity occurs.
\end{proof}

\begin{remark}
This reflects the fact that group directions associated with vanishing coordinates are suppressed by the metric.
\end{remark}

\subsection{Energy formulation}

\begin{definition}
For a tangent vector $(v_i,\xi_i)$, define the energy
\[
E(v,\xi)
=
\sum_i v_i^2
+
\sum_i x_i^2 \|\xi_i\|_G^2.
\]
\end{definition}

\begin{proposition}
The metric $g$ is the infinitesimal variation of this energy functional.
\end{proposition}

\subsection{Relation with the stratification}

\begin{proposition}
The metric does not distinguish strata at first order.
\end{proposition}

\begin{proof}
The tangent directions corresponding to inactive indices $i \notin J$ appear in the first term $\sum v_i^2$, but the group components only appear when $x_i \neq 0$. This matches the tangent structure described in Section~5.
\end{proof}

\subsection{Invariance properties}

\begin{proposition}
The metric $g$ is invariant under the action of $G$.
\end{proposition}

\begin{proof}
The action of $G$ only affects the $g_i$ coordinates, and the metric on each copy of $G$ is right-invariant.
\end{proof}

\begin{proposition}
The metric $g$ is invariant under the $\mathbb{Z}_2$-action.
\end{proposition}

\begin{proof}
The transformation $x_i \mapsto -x_i$ leaves both $x_i^2$ and $dx_i^2$ invariant.
\end{proof}
\subsection{Non-degeneracy of the metric}

We analyze the degeneracy properties of the metric introduced above.

\begin{proposition}
Let $p \in E^{\mathrm{sph}}(G)$ have support $J$. Then the metric $g$ satisfies:
\begin{itemize}
\item it is positive definite on
\[
\left\{
(v_i)_{i\in\mathbb{N}}^{\mathrm{fin}}
\;\middle|\;
\sum_{j\in J} x_j v_j = 0
\right\}
\oplus
\bigoplus_{j\in J} T_{g_j}G,
\]
\item it vanishes on directions corresponding to variations of $g_i$ for indices $i \notin J$.
\end{itemize}
\end{proposition}

\begin{proof}
By definition, the metric reads
\[
g(v,\xi)
=
\sum_i v_i^2
+
\sum_i x_i^2 \|\xi_i\|_G^2.
\]

If $i \notin J$, then $x_i = 0$, hence the corresponding term
\[
x_i^2 \|\xi_i\|_G^2
\]
vanishes.

On the other hand, for $j \in J$, one has $x_j \neq 0$, hence the contribution
\[
x_j^2 \|\xi_j\|_G^2
\]
is strictly positive unless $\xi_j = 0$.

The spherical part $\sum v_i^2$ is positive definite under the constraint
\[
\sum_{j\in J} x_j v_j = 0.
\]

This proves the result.
\end{proof}

\begin{corollary}
The metric $g$ is a pseudo-metric in the sense of diffeological vector pseudo-bundles.
\end{corollary}

\begin{remark}
The degeneracy directions correspond exactly to group components associated with vanishing spherical coordinates. These directions are invisible in the quotient defining $E^{\mathrm{sph}}(G)$.
\end{remark}

\begin{proposition}
The metric is non-degenerate when restricted to the tangent space of any stratum $E_J^\circ$.
\end{proposition}

\begin{proof}
On $E_J^\circ$, only the coordinates indexed by $J$ are active. The metric reduces to
\[
\sum_{j\in J} v_j^2
+
\sum_{j\in J} x_j^2 \|\xi_j\|_G^2,
\]
which is strictly positive definite.
\end{proof}

\begin{remark}
Thus, the degeneracy of the metric is purely transverse to the strata.
\end{remark}

\begin{proposition}
The kernel of $g$ at a point $p$ is given by
\[
\ker g_p
=
\bigoplus_{i\notin J} T_{g_i}G.
\]
\end{proposition}

\begin{proof}
Immediate from the expression of the metric.
\end{proof}

\begin{remark}
This kernel corresponds precisely to directions that are not detected at the level of the quotient structure.
\end{remark}
\subsection{Diffeological interpretation of non-degeneracy}

We reinterpret the degeneracy of the metric in terms of diffeological vector pseudo-bundles.

\subsubsection{Tangent pseudo-bundle}

Let
\[
T E^{\mathrm{sph}}(G) \longrightarrow E^{\mathrm{sph}}(G)
\]
denote the internal tangent pseudo-bundle.

\begin{proposition}
The assignment
\[
p \longmapsto \ker g_p
\]
defines a diffeological vector pseudo-bundle
\[
\mathcal{K} \subset T E^{\mathrm{sph}}(G).
\]
\end{proposition}

\begin{proof}
By Section~5, the tangent pseudo-bundle is locally described by
\[
T E^{\mathrm{sph}}(G) \simeq \varinjlim_I T(S_I \times G^I).
\]

On each local model, the kernel of $g$ is given by
\[
\bigoplus_{i \notin I_{\mathrm{active}}} T_{g_i}G,
\]
which varies smoothly with respect to the plots.

Compatibility under inclusions $J \subset I$ ensures that these local kernels glue into a well-defined pseudo-bundle.
\end{proof}

\subsubsection{Quotient pseudo-bundle}

\begin{definition}
Define the quotient pseudo-bundle
\[
\mathcal{H}
:=
T E^{\mathrm{sph}}(G) / \mathcal{K}.
\]
\end{definition}

\begin{proposition}
The metric $g$ induces a non-degenerate bilinear form
\[
\bar g : \mathcal{H} \times \mathcal{H} \longrightarrow \mathbb{R}.
\]
\end{proposition}

\begin{proof}
By construction, $g$ vanishes precisely on $\mathcal{K}$. Therefore it descends to the quotient.

Non-degeneracy follows from the positivity on active spherical and group directions.
\end{proof}

\subsubsection{Intrinsic characterization}

\begin{proposition}
The pseudo-bundle $\mathcal{K}$ coincides with the kernel of the differential of the projection
\[
\pi : \widetilde E^{\mathrm{sph}}(G) \to E^{\mathrm{sph}}(G).
\]
\end{proposition}

\begin{proof}
Directions corresponding to variations of $g_i$ with $x_i=0$ are precisely those that are collapsed by the quotient relation defining $E^{\mathrm{sph}}(G)$.

Hence they are tangent to the fibers of $\pi$, and therefore lie in the kernel of $d\pi$.
\end{proof}

\begin{corollary}
We have a short exact sequence of pseudo-bundles
\[
0 \longrightarrow \mathcal{K}
\longrightarrow T E^{\mathrm{sph}}(G)
\longrightarrow \mathcal{H}
\longrightarrow 0.
\]
\end{corollary}

\subsubsection{Geometric interpretation}

\begin{remark}
The pseudo-bundle $\mathcal{K}$ encodes directions that are invisible in the quotient defining $E^{\mathrm{sph}}(G)$.

The quotient $\mathcal{H}$ corresponds to the effective geometric directions on which the metric is non-degenerate.
\end{remark}

\begin{remark}
This decomposition is canonical and compatible with the diffeological structure, providing an intrinsic replacement for the usual splitting into tangent and normal directions.
\end{remark}
\subsection{Summary}

The spherical Milnor space carries a natural Riemannian structure defined by:
\begin{itemize}
\item the Euclidean metric on the spherical coordinates,
\item a weighted metric on the group components,
\item compatibility with the quotient and gluing structure,
\item smooth behavior across strata.
\end{itemize}

\begin{remark}
This metric provides the geometric foundation for the differential operators constructed in the following sections.
\end{remark}
\section{Differential forms and Hodge-type structures} \label{7}

\subsection{Forms on a diffeological space}

Let $X$ be a diffeological space. Recall that a differential $k$-form on $X$ is given by an assignment
\[
p : U \to X \quad \longmapsto \quad \omega_p \in \Omega^k(U)
\]
compatible with smooth reparametrizations.

We denote by
\[
\Omega^k(X)
\]
the space of smooth $k$-forms on $X$.

\subsection{Horizontal forms}

Let
\[
\mathcal K \subset TE^{\mathrm{sph}}(G)
\]
be the degeneracy pseudo-bundle defined in Section~6.

\begin{definition}
A differential form $\omega \in \Omega^k(E^{\mathrm{sph}}(G))$ is called \emph{horizontal} if
\[
\iota_v \omega = 0
\quad \text{for all } v \in \mathcal K.
\]
\end{definition}

We denote by
\[
\Omega^k_{\mathrm{hor}}(E^{\mathrm{sph}}(G))
\]
the space of horizontal forms.

\begin{proposition}
The space $\Omega^\bullet_{\mathrm{hor}}(E^{\mathrm{sph}}(G))$ is stable under the exterior differential $d$.
\end{proposition}

\begin{proof}
If $\omega$ vanishes on $\mathcal K$, then so does $d\omega$ by standard Cartan calculus.
\end{proof}

\subsection{Metric pairing on horizontal forms}

Let
\[
\mathcal H
=
TE^{\mathrm{sph}}(G)/\mathcal K
\]
be the effective horizontal pseudo-bundle endowed with the non-degenerate
metric
\[
\bar g.
\]

\subsubsection{Musical map associated with the metric}

The metric defines the fiberwise bundle morphism
\[
\flat_{\bar g}:
\mathcal H
\longrightarrow
\mathcal H^*,
\qquad
v\mapsto \bar g(v,\cdot).
\]

\begin{definition}
We say that the metric \(\bar g\) is:
\begin{itemize}
\item weakly non-degenerate if
\[
\flat_{\bar g}
\]
is injective;
\item strongly non-degenerate if
\[
\flat_{\bar g}
\]
is an isomorphism onto the dual pseudo-bundle.
\end{itemize}
\end{definition}

\begin{remark}
In finite-dimensional local models, weak and strong non-degeneracy coincide.
In infinite-dimensional settings, the distinction becomes essential.
\end{remark}

\subsubsection{Pairing on horizontal one-forms}

Assume from now on that \(\bar g\) is strongly non-degenerate.

\begin{definition}
The inverse musical morphism
\[
\sharp_{\bar g}:
\mathcal H^*
\longrightarrow
\mathcal H
\]
is defined as the inverse of
\[
\flat_{\bar g}.
\]

For horizontal one-forms
\[
\alpha,\beta\in\Gamma(\mathcal H^*),
\]
define the metric pairing by
\[
\langle \alpha,\beta\rangle_{\bar g}
:=
\bar g(\sharp_{\bar g}\alpha,\sharp_{\bar g}\beta).
\]
\end{definition}

\begin{proposition}
The pairing
\[
\langle\cdot,\cdot\rangle_{\bar g}
\]
is symmetric and non-degenerate.
\end{proposition}

\begin{proof}
Symmetry follows from symmetry of the metric \(\bar g\). Non-degeneracy follows
from invertibility of the musical isomorphism
\[
\sharp_{\bar g}.
\]
\end{proof}

\begin{remark}
This formulation avoids introducing a formal inverse tensor
\[
\bar g^{-1},
\]
which may not exist globally in the diffeological or infinite-dimensional
setting.
\end{remark}

\subsubsection{Extension to higher-degree forms}

The metric pairing extends fiberwise to exterior powers
\[
\Lambda^k\mathcal H^*
\]
by multilinearity and determinant contraction.

Explicitly, if
\[
\alpha=\alpha_1\wedge\cdots\wedge\alpha_k,
\qquad
\beta=\beta_1\wedge\cdots\wedge\beta_k,
\]
define
\[
\langle\alpha,\beta\rangle_{\bar g}
=
\det\!\Big(
\langle\alpha_i,\beta_j\rangle_{\bar g}
\Big)_{1\leq i,j\leq k}.
\]

\subsection{Formal adjoint of the exterior differential}

We now define the codifferential without using a Hodge star operator.

\subsubsection{Finite-rank case}

Assume first that the horizontal geometry is locally of finite rank and that a
fiberwise density or volume form has been fixed.

\begin{definition}
Let
\[
d:
\Omega^k_{\mathrm{hor}}
(E^{\mathrm{sph}}(G))
\longrightarrow
\Omega^{k+1}_{\mathrm{hor}}
(E^{\mathrm{sph}}(G))
\]
be the horizontal exterior differential.

The codifferential
\[
\delta:
\Omega^{k+1}_{\mathrm{hor}}
(E^{\mathrm{sph}}(G))
\longrightarrow
\Omega^{k}_{\mathrm{hor}}
(E^{\mathrm{sph}}(G))
\]
is defined as the formal adjoint of \(d\) with respect to the metric pairing
and the chosen density.
\end{definition}

\begin{proposition}
In the finite-rank setting, the codifferential \(\delta\) is well-defined.
\end{proposition}

\begin{proof}
On each finite local model, the horizontal geometry is an ordinary
finite-dimensional Riemannian geometry. The metric pairing together with the
chosen density defines an \(L^2\)-type pairing on horizontal forms.

The formal adjoint of \(d\) therefore exists locally by the standard
integration-by-parts construction. Compatibility of the local structures under
the transition maps implies that these local codifferentials glue.
\end{proof}

\begin{remark}
The codifferential depends not only on the metric but also on the choice of
density or integration functional used to define the formal adjoint.
\end{remark}

\subsubsection{Infinite-dimensional case}

In infinite dimension, the existence of a formal adjoint is substantially more
delicate.

\begin{remark}
Even if the metric is strongly non-degenerate, a codifferential operator does
not automatically exist in the absence of:
\begin{itemize}
\item a suitable measure or density theory;
\item an \(L^2\)-type completion;
\item and adequate functional-analytic properties of the differential complex.
\end{itemize}
\end{remark}

\begin{definition}
An admissible Hodge datum on the horizontal pseudo-bundle consists of:
\begin{itemize}
\item a strongly non-degenerate metric;
\item a compatible density or trace functional;
\item a Hilbert-type completion of horizontal forms;
\item and a densely defined differential operator \(d\).
\end{itemize}

Under such assumptions, the codifferential is defined as the Hilbert adjoint
of \(d\).
\end{definition}

\begin{remark}
Thus the codifferential should be regarded as an additional analytic structure,
rather than as a purely formal consequence of the metric.
\end{remark}
\subsection{Hodge-type Laplacian}

We now define a Laplace-type operator associated with the horizontal
differential complex.

\subsubsection{Finite-rank case}

Assume first that:
\begin{itemize}
\item the effective horizontal pseudo-bundle
\[
\mathcal H
=
TE^{\mathrm{sph}}(G)/\mathcal K
\]
has finite rank on the considered local models;
\item the metric \(\bar g\) is strongly non-degenerate;
\item a compatible density or integration functional has been fixed;
\item the codifferential
\[
\delta
\]
is defined as the formal adjoint of
\[
d.
\]
\end{itemize}

\begin{definition}
The horizontal Hodge-type Laplacian is
\[
\Delta
:=
d\delta+\delta d
\]
acting on horizontal differential forms.
\end{definition}

\begin{proposition}
Under the preceding assumptions, the operator
\[
\Delta
=
d\delta+\delta d
\]
is well-defined and formally self-adjoint.
\end{proposition}

\begin{proof}
By construction, \(\delta\) is the formal adjoint of \(d\) with respect to
the chosen pairing on horizontal forms. Therefore:
\[
\langle d\alpha,\beta\rangle
=
\langle \alpha,\delta\beta\rangle
\]
for admissible forms.

Hence
\[
(d\delta)^*=\delta^*d^*=d\delta
\]
and similarly
\[
(\delta d)^*=d^*\delta^*=\delta d.
\]

Therefore
\[
\Delta^*
=
(d\delta+\delta d)^*
=
d\delta+\delta d
=
\Delta.
\]
Thus \(\Delta\) is formally self-adjoint.
\end{proof}

\subsection{Local expression}

We now describe the local structure of the Laplacian.

\begin{proposition}
On a finite local model
\[
S_I\times G^I,
\]
equipped with the warped metric
\[
g
=
g_{S_I}
+
\sum_{i\in I}x_i^2g_G^{(i)},
\]
the Hodge-type Laplacian decomposes locally as
\[
\Delta
=
\Delta_{S_I}
+
\sum_{i\in I}\frac1{x_i^2}\Delta_{G,i}
+
L,
\]
where:
\begin{itemize}
\item \(\Delta_{S_I}\) is the Laplacian associated with the spherical part;
\item \(\Delta_{G,i}\) is the Laplacian associated with the \(i\)-th copy of
\(G\);
\item \(L\) is a lower-order differential operator involving derivatives of
the barycentric coefficients \(x_i\).
\end{itemize}
\end{proposition}

\begin{proof}
The inverse metric associated with
\[
g
=
g_{S_I}
+
\sum_i x_i^2g_G^{(i)}
\]
takes the form
\[
g^{-1}
=
g_{S_I}^{-1}
+
\sum_i \frac1{x_i^2}(g_G^{(i)})^{-1}.
\]

Consequently, the second-order contribution of the Laplacian in the group
directions carries the factors
\[
\frac1{x_i^2}.
\]

The derivatives of the coefficients \(x_i^2\) contribute additional first-
and zeroth-order terms, which are collected into the operator \(L\).
\end{proof}

\begin{remark}
The factors
\[
\frac1{x_i^2}
\]
are characteristic of warped-product geometries.
\end{remark}

\begin{remark}
The local formula should be understood at the level of finite local models.
In infinite-dimensional settings, additional analytic assumptions are required
in order to define the Laplacian rigorously.
\end{remark}

\subsection{Behavior near boundary strata}

The coefficients
\[
\frac1{x_i^2}
\]
appear singular near the loci where
\[
x_i=0.
\]

However, these apparent singularities are compensated by the degeneration of
the corresponding horizontal directions.

\begin{proposition}
The horizontal Hodge-type Laplacian extends through the boundary strata of the
spherical Milnor model.
\end{proposition}

\begin{proof}
When
\[
x_i\to0,
\]
the metric contribution
\[
x_i^2g_G^{(i)}
\]
collapses the corresponding horizontal directions.

Thus the singular coefficient
\[
\frac1{x_i^2}
\]
acts only on directions whose effective geometric contribution vanishes in the
quotient horizontal geometry.

Equivalently, the degeneration of the metric compensates the blow-up of the
inverse metric coefficients, so that the resulting operator remains compatible
with the effective horizontal pseudo-bundle.
\end{proof}

\begin{remark}
This cancellation mechanism is analogous to standard phenomena in conic and
warped geometries.
\end{remark}

\subsection{Harmonic horizontal forms}

\begin{definition}
A horizontal differential form
\[
\omega
\]
is called harmonic if
\[
\Delta\omega=0.
\]
\end{definition}

\begin{remark}
The notion of harmonicity depends only on the effective horizontal geometry
encoded by
\[
\mathcal H
=
TE^{\mathrm{sph}}(G)/\mathcal K.
\]
\end{remark}

\begin{remark}
In infinite-dimensional settings, the analytic meaning of harmonicity depends
on the chosen Hilbert completion and on the domain of the Laplacian.
\end{remark}

\subsection{Interpretation}

\begin{remark}
The preceding construction provides a Hodge-type theory without introducing a
Hodge star operator.
\end{remark}

\begin{remark}
This is particularly useful in diffeological and infinite-dimensional
settings, where:
\begin{itemize}
\item duality operators may fail to exist globally;
\item top-degree forms may not be available;
\item or the tangent geometry may have non-constant rank.
\end{itemize}
\end{remark}

\begin{remark}
The restriction to horizontal forms removes the degenerate directions
associated with the kernel pseudo-bundle
\[
\mathcal K,
\]
which ensures that the effective metric geometry remains non-degenerate.
\end{remark}

\subsection{Summary}

Under the finite-rank or admissible analytic assumptions discussed above, we
have constructed:
\begin{itemize}
\item a horizontal differential complex;
\item a metric pairing on horizontal forms;
\item a codifferential defined as a formal adjoint;
\item a Hodge-type Laplacian adapted to the horizontal geometry;
\item and a corresponding notion of harmonic horizontal forms.
\end{itemize}


These constructions provide the analytic framework required for the
Dirac-type and spinorial operators introduced in the next section.
Therefore, all along next section we assume that Hodge structures and related operators are well defined. 
\section{Dirac-type operators on the spherical Milnor space} \label{8}

\subsection{Clifford structure on the horizontal pseudo-bundle}

Let
\[
\mathcal H = TE^{\mathrm{sph}}(G)/\mathcal K
\]
be the horizontal pseudo-bundle endowed with the non-degenerate metric $\bar g$.

\begin{definition}
The Clifford algebra bundle $\mathrm{Cl}(\mathcal H)$ is defined fiberwise by the relations
\[
v \cdot w + w \cdot v = -2 \bar g(v,w),
\qquad v,w \in \mathcal H.
\]
\end{definition}

\begin{remark}
This construction is well-defined since $\bar g$ is non-degenerate on $\mathcal H$.
\end{remark}

\subsection{Clifford modules}

\begin{definition}
A \emph{Clifford module} over $E^{\mathrm{sph}}(G)$ is a diffeological vector pseudo-bundle
\[
\mathcal E \longrightarrow E^{\mathrm{sph}}(G)
\]
equipped with a smooth fiberwise action of $\mathrm{Cl}(\mathcal H)$.
\end{definition}

\begin{remark}
In local charts $S_I \times G^I$, one can take
\[
\mathcal E = \Lambda^\bullet \mathcal H^*
\]
or spinor-type constructions when available.
\end{remark}

\subsection{Connection on the horizontal bundle}

Let us first give a statement applicable when $G$ is a classical, finite dimensional Lie group.

\begin{proposition}
Let \(\mathcal H\to E^{\mathrm{sph}}(G)\) be the effective horizontal
pseudo-bundle endowed with the non-degenerate metric \(\bar g\). Assume that:
\begin{enumerate}
\item on each finite local model \(S_I\times G^I\), the bundle
\(\mathcal H\) is represented by a finite-dimensional Riemannian vector bundle;
\item the transition maps preserve \(\bar g\);
\item \(\mathcal E\to E^{\mathrm{sph}}(G)\) is a Clifford module associated
with \(\mathrm{Cl}(\mathcal H)\);
\item the local Clifford modules glue compatibly with the transition maps.
\end{enumerate}
Then there exists a connection
\[
\nabla^{\mathcal E}:\Gamma(\mathcal E)
\longrightarrow
\Gamma(\mathcal H^*\otimes \mathcal E)
\]
compatible with the Clifford structure, i.e.
\[
\nabla^{\mathcal E}_X(c(v)s)
=
c(\nabla^{\mathcal H}_X v)s
+
c(v)\nabla^{\mathcal E}_X s
\]
for all local sections \(X,v\) of \(\mathcal H\) and all local sections
\(s\) of \(\mathcal E\).
\end{proposition}

\begin{proof}
We proceed locally and then check compatibility under gluing.

Let \(S_I\times G^I\) be one of the finite-dimensional local models of
\(E^{\mathrm{sph}}(G)\). On this model, the horizontal pseudo-bundle
\(\mathcal H\) is represented by a finite-dimensional vector bundle endowed
with the non-degenerate metric \(\bar g\). Therefore, by the ordinary
fundamental theorem of Riemannian geometry, there exists a unique
Levi-Civita connection
\[
\nabla^{\mathcal H,I}
\]
on \(\mathcal H|_{S_I\times G^I}\), characterized by
\[
\nabla^{\mathcal H,I}\bar g=0
\]
and by the vanishing of torsion.

This connection induces canonically a connection on the Clifford algebra
bundle
\[
\mathrm{Cl}(\mathcal H)|_{S_I\times G^I}.
\]
Indeed, for a local section \(v\) of \(\mathcal H\), one defines
\[
\nabla^{\mathrm{Cl},I}_X c(v)
=
c(\nabla^{\mathcal H,I}_X v),
\]
and extends this rule to products by the Leibniz rule. This is well-defined
because \(\nabla^{\mathcal H,I}\) preserves the metric. To see this, recall
that the Clifford relation is
\[
c(v)c(w)+c(w)c(v)=-2\bar g(v,w).
\]
Applying \(\nabla^{\mathrm{Cl},I}_X\) gives
\[
c(\nabla_X v)c(w)+c(v)c(\nabla_X w)
+
c(\nabla_X w)c(v)+c(w)c(\nabla_X v)
=
-2X(\bar g(v,w)).
\]
Since \(\nabla^{\mathcal H,I}\bar g=0\), the right-hand side equals
\[
-2\bar g(\nabla_X v,w)-2\bar g(v,\nabla_X w),
\]
which is exactly the derivative of the Clifford relation. Hence the
connection preserves the Clifford algebra structure.

Now let \(\mathcal E|_{S_I\times G^I}\) be the corresponding local
Clifford module. By assumption, \(\mathcal E\) is associated to the
Clifford structure. Locally, one may choose a connection
\[
\nabla^{\mathcal E,I}
\]
on \(\mathcal E|_{S_I\times G^I}\) compatible with
\(\nabla^{\mathrm{Cl},I}\), meaning that
\[
\nabla^{\mathcal E,I}_X(c(a)s)
=
c(\nabla^{\mathrm{Cl},I}_X a)s
+
c(a)\nabla^{\mathcal E,I}_X s
\]
for every local Clifford algebra section \(a\) and every local module
section \(s\). In particular, for \(a=v\in\Gamma(\mathcal H)\), this gives
\[
\nabla^{\mathcal E,I}_X(c(v)s)
=
c(\nabla^{\mathcal H,I}_X v)s
+
c(v)\nabla^{\mathcal E,I}_X s.
\]

It remains to check that these local connections glue. Let
\(S_J\times G^J\) be another local model and suppose the two charts overlap
through a refinement or inclusion of supports. By hypothesis, the transition
maps preserve the metric \(\bar g\), hence they identify the corresponding
Levi-Civita connections:
\[
\nabla^{\mathcal H,I}
=
\nabla^{\mathcal H,J}
\]
on the overlap, after transport by the transition map. Consequently the
induced Clifford algebra connections also agree on overlaps.

Since the Clifford modules are assumed to glue compatibly with the same
transition maps, the local module connections
\(\nabla^{\mathcal E,I}\) and \(\nabla^{\mathcal E,J}\) agree on overlaps.
Therefore they define a global connection
\[
\nabla^{\mathcal E}:\Gamma(\mathcal E)
\to
\Gamma(\mathcal H^*\otimes\mathcal E).
\]

The compatibility identity is local, and it has been verified on every
finite-dimensional local model. Since the diffeology on
\(E^{\mathrm{sph}}(G)\) is generated by these local models, the identity
holds globally. Hence \(\nabla^{\mathcal E}\) is a Clifford-compatible
connection.
\end{proof}

\begin{remark}
The proposition should not be read as asserting that every abstract
Clifford module over an arbitrary diffeological pseudo-bundle admits a
compatible connection. The result applies to the Clifford modules obtained
from the local finite-dimensional models used in the construction, provided
that the metric, Clifford action, and transition data are compatible.
\end{remark}
Let us now precise the statements for more general Lie groups, in particular infinite dimensional ones. 
\begin{proposition}
Let \(G\) be a diffeological Lie group, possibly infinite-dimensional.
Assume that the effective horizontal pseudo-bundle
\[
\mathcal H \to E^{\mathrm{sph}}(G)
\]
is endowed with a non-degenerate metric \(\bar g\), and assume moreover that
there exists a metric connection
\[
\nabla^{\mathcal H}:\Gamma(\mathcal H)\to
\Gamma(\mathcal H^*\otimes \mathcal H)
\]
satisfying
\[
\nabla^{\mathcal H}\bar g=0.
\]
Let \(\mathcal E\to E^{\mathrm{sph}}(G)\) be a Clifford module associated with
\(\mathrm{Cl}(\mathcal H)\), and assume that \(\mathcal E\) is equipped locally
with transition data compatible with \(\nabla^{\mathcal H}\).

Then there exists a Clifford-compatible connection
\[
\nabla^{\mathcal E}:\Gamma(\mathcal E)\to
\Gamma(\mathcal H^*\otimes \mathcal E)
\]
satisfying
\[
\nabla^{\mathcal E}_X(c(v)s)
=
c(\nabla^{\mathcal H}_Xv)s
+
c(v)\nabla^{\mathcal E}_Xs.
\]
\end{proposition}
\begin{proof}
The key point is that, in the infinite-dimensional case, the existence of a
Levi-Civita connection is not automatic. We therefore assume the existence of
a metric connection \(\nabla^{\mathcal H}\) on \(\mathcal H\).

Since \(\nabla^{\mathcal H}\bar g=0\), it induces a connection on the Clifford
algebra bundle \(\mathrm{Cl}(\mathcal H)\). Indeed, for a section
\(v\in\Gamma(\mathcal H)\), define
\[
\nabla^{\mathrm{Cl}}_X c(v):=
c(\nabla^{\mathcal H}_Xv),
\]
and extend this rule to products by the Leibniz rule.

This is compatible with the Clifford relation
\[
c(v)c(w)+c(w)c(v)=-2\bar g(v,w),
\]
because metric compatibility gives
\[
X\bar g(v,w)
=
\bar g(\nabla^{\mathcal H}_Xv,w)
+
\bar g(v,\nabla^{\mathcal H}_Xw).
\]
Hence the induced connection preserves the Clifford algebra structure.

By assumption, the Clifford module \(\mathcal E\) is locally compatible with
this Clifford algebra connection. Therefore one can choose local module
connections \(\nabla^{\mathcal E}\) satisfying
\[
\nabla^{\mathcal E}_X(c(a)s)
=
c(\nabla^{\mathrm{Cl}}_Xa)s
+
c(a)\nabla^{\mathcal E}_Xs.
\]
For \(a=v\in\Gamma(\mathcal H)\), this gives the desired formula.

The transition data are assumed compatible with \(\nabla^{\mathcal H}\) and
with the Clifford action; hence the local module connections glue. This
defines a global Clifford-compatible connection on \(\mathcal E\).
\end{proof}
\subsection{Definition of the Dirac operator}

\subsection{Dirac-type operators}

We now define Dirac-type operators associated with the horizontal pseudo-bundle
\[
\mathcal H\to E^{\mathrm{sph}}(G).
\]

The construction must be treated with some care. If \(\mathcal H\) has finite
rank on each local model, the usual Clifford contraction gives a well-defined
Dirac operator. In the infinite-dimensional case, additional analytic
assumptions are required in order to make sense of the infinite Clifford
contraction.

\subsubsection{The finite-rank case}

Assume first that \(\mathcal H\) is locally of finite rank and is equipped with:
\begin{itemize}
\item a non-degenerate metric \(\bar g\),
\item a Clifford module \(\mathcal E\),
\item a Clifford-compatible connection \(\nabla^{\mathcal E}\).
\end{itemize}

\begin{definition}
Let \((e_1,\ldots,e_r)\) be a local \(\bar g\)-orthonormal frame of
\(\mathcal H\). The associated Dirac operator is
\[
D
=
\sum_{a=1}^r c(e_a)\nabla^{\mathcal E}_{e_a},
\]
where \(c(e_a)\) denotes Clifford multiplication.
\end{definition}

\begin{proposition}
The operator
\[
D=\sum_{a=1}^r c(e_a)\nabla^{\mathcal E}_{e_a}
\]
is independent of the choice of local orthonormal frame.
\end{proposition}

\begin{proof}
Let \((e_a)\) and \((e'_a)\) be two local orthonormal frames. Then there exists
a smooth map
\[
A=(A_{ab})
\]
with values in the orthogonal group \(O(r)\) such that
\[
e'_a=\sum_b A_{ab}e_b.
\]
Since the connection is linear in the vector-field argument,
\[
\nabla^{\mathcal E}_{e'_a}
=
\sum_c A_{ac}\nabla^{\mathcal E}_{e_c}.
\]
Similarly, Clifford multiplication is linear:
\[
c(e'_a)
=
\sum_b A_{ab}c(e_b).
\]
Therefore
\[
\sum_a c(e'_a)\nabla^{\mathcal E}_{e'_a}
=
\sum_{a,b,c}A_{ab}A_{ac}c(e_b)\nabla^{\mathcal E}_{e_c}.
\]
Since \(A\in O(r)\),
\[
\sum_a A_{ab}A_{ac}=\delta_{bc}.
\]
Hence
\[
\sum_a c(e'_a)\nabla^{\mathcal E}_{e'_a}
=
\sum_b c(e_b)\nabla^{\mathcal E}_{e_b}.
\]
Thus the local expression for \(D\) is frame-independent.
\end{proof}

\begin{proposition}
The local operators glue to a global operator
\[
D:\Gamma(\mathcal E)\longrightarrow \Gamma(\mathcal E)
\]
whenever the metric, Clifford module and connection glue compatibly with the
diffeological transition maps.
\end{proposition}

\begin{proof}
The definition of \(D\) is local, and the previous proposition shows that it
does not depend on the choice of orthonormal frame. On overlaps of local
models, the metric, Clifford multiplication and connection are assumed to be
identified by the transition maps. Hence the local expressions agree on
overlaps. Therefore they glue to a globally defined operator.
\end{proof}

\subsubsection{The infinite-dimensional case}

When \(G\) is infinite-dimensional, the horizontal pseudo-bundle
\(\mathcal H\) may be infinite-dimensional. In that case, the formal expression
\[
\sum_a c(e_a)\nabla^{\mathcal E}_{e_a}
\]
is generally not meaningful without additional convergence or trace-class
assumptions.

\begin{definition}
An infinite-dimensional Clifford-compatible Dirac datum consists of:
\begin{itemize}
\item a Hilbert or convenient completion \(\widehat{\mathcal H}\) of
\(\mathcal H\);
\item a non-degenerate metric or inner product on \(\widehat{\mathcal H}\);
\item a Clifford module \(\mathcal E\);
\item a compatible connection \(\nabla^{\mathcal E}\);
\item a choice of admissible orthonormal frame \((e_a)_{a\geq1}\);
\item a convergence condition ensuring that
\[
\sum_{a\geq1} c(e_a)\nabla^{\mathcal E}_{e_a}
\]
converges in a specified operator topology.
\end{itemize}
\end{definition}

\begin{definition}
Under such admissibility assumptions, the infinite-dimensional Dirac-type
operator is defined by
\[
D
=
\overline{\sum_{a\geq1} c(e_a)\nabla^{\mathcal E}_{e_a}},
\]
where the bar denotes the closure of the densely defined operator obtained
from the convergent series.
\end{definition}

\begin{remark}
In the infinite-dimensional case, frame independence is not automatic. It must
be required that changes of admissible frames act by transformations preserving
the convergence class of the series, for example by a restricted orthogonal
group or another suitable subgroup of the full orthogonal group.
\end{remark}

\begin{proposition}
Assume that admissible changes of frame belong to a restricted orthogonal group
under which the Clifford contraction series is invariant. Then the operator
\(D\) is independent of the admissible frame.
\end{proposition}

\begin{proof}
The proof is formally the same as in finite dimension, but one must justify
interchanging infinite sums. This is precisely the role of the restricted
orthogonal admissibility assumption. It ensures that the matrix coefficients
of a change of frame preserve the convergence of the series and that the
orthogonality identity
\[
\sum_a A_{ab}A_{ac}=\delta_{bc}
\]
is valid in the chosen topology. Hence the transformed Clifford contraction
equals the original one.
\end{proof}

\begin{remark}
Thus, in the infinite-dimensional setting, the Dirac operator is not a purely
formal consequence of the metric and Clifford structure. It requires an
additional analytic datum specifying the class of admissible frames and the
operator topology in which the contraction converges.
\end{remark}

\subsubsection{Dirac operators on finite local models}

In the present article, the primary construction is performed on the finite
local models
\[
S_I\times G^I.
\]
If \(G\) is finite-dimensional, these are finite-dimensional manifolds or
finite-dimensional diffeological models, and the preceding finite-rank
construction applies directly.

If \(G\) is infinite-dimensional, the local models remain finite in the Milnor
index set \(I\), but the factors \(G^I\) may still be infinite-dimensional.
Therefore the Dirac construction requires either:
\begin{itemize}
\item an explicit Hilbert completion of the tangent spaces of \(G\);
\item a compatible metric connection;
\item a Clifford module;
\item and an admissibility condition for the Clifford contraction.
\end{itemize}

\textbf{conclusion}
For finite-dimensional structure groups, the Dirac operator is canonically
defined from the metric, Clifford module and compatible connection. For
infinite-dimensional structure groups, the same notation \(D\) will be used
only after an admissible infinite-dimensional Dirac datum has been fixed.


\subsection{Local expression of the Dirac operator}

Let
\[
S_I\times G^I
\]
be a finite local model of the spherical Milnor space. We denote by
\[
(x_i)_{i\in I}
\]
the spherical barycentric coordinates and by
\[
\mathfrak g
\]
the Lie algebra of \(G\).

Assume that:
\begin{itemize}
\item \(G\) is equipped with a left-invariant metric;
\item the horizontal pseudo-bundle \(\mathcal H\) is endowed with the
barycentric metric
\[
\bar g
=
g_{S_I}
+
\sum_{i\in I} x_i^2\, g_G^{(i)};
\]
\item a Clifford-compatible connection has been fixed.
\end{itemize}

\begin{proposition}
Locally, the Dirac operator decomposes as
\[
D
=
D_{S_I}
+
\sum_{i\in I}\frac{1}{x_i}D_{G,i}
+
L,
\]
where:
\begin{itemize}
\item \(D_{S_I}\) is the Dirac operator associated with the spherical part;
\item \(D_{G,i}\) is the Dirac operator acting on the \(i\)-th copy of \(G\);
\item \(L\) is a first-order operator involving derivatives of the barycentric
coefficients \(x_i\).
\end{itemize}
\end{proposition}

\begin{proof}
Let
\[
(e_\alpha)
\]
be a local orthonormal frame for the spherical factor \(S_I\), and let
\[
(f_a^{(i)})
\]
be orthonormal frames for the metric \(g_G^{(i)}\) on the \(i\)-th copy of
\(G\).

Because the metric on the \(i\)-th group factor is multiplied by \(x_i^2\),
the corresponding orthonormal frame for the total metric \(\bar g\) is
\[
\widetilde f_a^{(i)}
=
\frac{1}{x_i}f_a^{(i)}
\]
whenever \(x_i\neq0\).

Hence the local Clifford contraction gives
\[
D
=
\sum_\alpha c(e_\alpha)\nabla_{e_\alpha}
+
\sum_{i,a}
c\!\left(\frac1{x_i}f_a^{(i)}\right)
\nabla_{\frac1{x_i}f_a^{(i)}}.
\]

Using linearity of Clifford multiplication and of the connection in the vector
field argument,
\[
D
=
D_{S_I}
+
\sum_i \frac1{x_i}D_{G,i}
+
L,
\]
where \(L\) collects the additional first-order terms produced by derivatives
of the scaling coefficients \(x_i\) through the connection.
\end{proof}

\begin{remark}
The factors \(1/x_i\) arise from the rescaling of the orthonormal frame under
the warped product-type metric
\[
g_{S_I}+\sum_i x_i^2 g_G^{(i)}.
\]
Thus they are geometric rather than singularity-theoretic in origin.
\end{remark}

\subsection{Square of the Dirac operator}

The classical Lichnerowicz formula does not automatically extend to arbitrary
diffeological or infinite-dimensional settings. We therefore restrict
ourselves to the finite-rank case.

\begin{proposition}
Assume that the horizontal pseudo-bundle \(\mathcal H\) has finite rank on the
considered local model and that the Clifford-compatible connection is induced
by a metric connection. Then locally,
\[
D^2
=
\nabla^*\nabla
+
\mathcal R,
\]
where:
\begin{itemize}
\item \(\nabla^*\nabla\) is the Bochner Laplacian associated with the chosen
connection;
\item \(\mathcal R\) is the curvature contribution induced by the Clifford
action of the curvature tensor.
\end{itemize}
\end{proposition}

\begin{proof}
On each finite local model, the metric and Clifford structures are classical.
The local Dirac operator is therefore an ordinary Dirac-type operator on a
finite-dimensional Riemannian manifold with warped product metric. The usual
Weitzenböck--Lichnerowicz argument applies and gives
\[
D^2
=
\nabla^*\nabla+\mathcal R.
\]
The precise expression of \(\mathcal R\) depends on the chosen Clifford module
and connection.
\end{proof}

\begin{remark}
In the spin case, the curvature contribution reduces locally to the usual
scalar-curvature term
\[
\frac14\operatorname{Scal}.
\]
However, we shall not require this more specific formula in the sequel.
\end{remark}

\begin{remark}
In infinite dimension, the expression \(D^2\) requires additional analytic
assumptions on domains, closures, and convergence of the Clifford contraction.
For this reason, the above formula should be understood as a local
finite-rank identity.
\end{remark}

\subsection{Behavior near the boundary strata}

We now analyze the behavior of the Dirac operator near loci where some
barycentric coordinates vanish.

\begin{proposition}
Let
\[
x_i=0
\]
for some index \(i\). Then the apparent singularity in the term
\[
\frac1{x_i}D_{G,i}
\]
is compensated by the degeneration of the corresponding horizontal directions.
Consequently, the induced Dirac operator on the effective horizontal
pseudo-bundle extends continuously across the boundary stratum.
\end{proposition}

\begin{proof}
By construction of the spherical Milnor quotient, when \(x_i=0\), the
corresponding group component becomes irrelevant in the equivalence relation.
Geometrically, the metric contribution
\[
x_i^2 g_G^{(i)}
\]
collapses the corresponding horizontal directions.

The vector fields
\[
\widetilde f_a^{(i)}
=
\frac1{x_i}f_a^{(i)}
\]
form an orthonormal frame only away from the locus \(x_i=0\). However, the
effective horizontal pseudo-bundle itself loses the corresponding directions
at the boundary.

Thus the apparently singular term acts on directions whose norm tends to zero
in the collapsing geometry. The resulting Clifford contraction therefore
extends through the quotient geometry of the horizontal pseudo-bundle.

Equivalently, the singular coefficient is compensated by the vanishing of the
metric weight defining the corresponding component of the horizontal geometry.
\end{proof}

\begin{remark}
The boundary behavior is analogous to the standard cancellation phenomena
occurring in warped product geometries and in conic degenerations.
\end{remark}

\begin{remark}
This cancellation mechanism is one of the main geometric advantages of the
spherical Milnor model. The degeneration of the barycentric metric is
compatible with the quotient structure and prevents the appearance of genuine
operator singularities on the effective horizontal geometry.
\end{remark}
\subsection{Invariance properties}

\begin{proposition}
The Dirac operator is invariant under the $G$-action.
\end{proposition}

\begin{proposition}
The Dirac operator is invariant under the $\mathbb{Z}_2$-action.
\end{proposition}

\subsection{Interpretation}

\begin{remark}
The operator $D$ acts only on the effective geometry encoded in $\mathcal H$, ignoring the degenerate directions.
\end{remark}

\begin{remark}
This construction provides a natural extension of Dirac operators to diffeological spaces with degenerate metrics.
\end{remark}
\subsection{Spin structures on the horizontal pseudo-bundle}

We now investigate the existence of spin structures on the effective
horizontal pseudo-bundle
\[
\mathcal H
=
TE^{\mathrm{sph}}(G)/\mathcal K.
\]

Since \(\mathcal H\) may be only a diffeological pseudo-bundle and may even be
infinite-dimensional, the construction must again be understood locally on the
finite Milnor models.

\subsubsection{Orthogonal structure}

Assume that the horizontal pseudo-bundle \(\mathcal H\) is endowed with a
non-degenerate metric
\[
\bar g.
\]

\begin{proposition}
On each finite local model
\[
S_I\times G^I,
\]
the metric \(\bar g\) defines a local orthonormal frame bundle.
\end{proposition}

\begin{proof}
On every finite local model, the horizontal pseudo-bundle is represented by a
finite-rank vector bundle equipped with the non-degenerate metric \(\bar g\).
Hence the usual orthonormal frame bundle construction applies locally.
\end{proof}

\begin{remark}
Globally, one should rather speak of a diffeological orthogonal
pseudo-bundle, since the rank of the horizontal geometry may vary across the
boundary strata.
\end{remark}

\subsubsection{Local structure of the horizontal pseudo-bundle}

\begin{proposition}
On a finite local model
\[
S_I\times G^I,
\]
the horizontal pseudo-bundle is locally represented by
\[
TS_I
\oplus
\bigoplus_{i\in I} TG,
\]
equipped with the warped metric
\[
g
=
g_{S_I}
\oplus
\bigoplus_{i\in I} x_i^2 g_G.
\]
\end{proposition}

\begin{proof}
This follows directly from the construction of the horizontal geometry and
from the decomposition of the barycentric metric into spherical and group
components.
\end{proof}

\begin{remark}
If \(G\) is finite-dimensional, then each local model is finite-dimensional
and classical spin geometry applies locally.
\end{remark}

\begin{remark}
If \(G\) is infinite-dimensional, the tangent bundle \(TG\) generally requires
additional Hilbert or convenient completions before Clifford and spin
structures can be defined analytically.
\end{remark}

\subsubsection{Local spin structures}

\begin{definition}
A local spin structure on a finite local model is a lift of the local
orthonormal frame bundle from
\[
O(r)
\]
to
\[
\operatorname{Spin}(r),
\]
where \(r\) denotes the local rank of the horizontal geometry.
\end{definition}

\begin{proposition}
Assume that:
\begin{itemize}
\item the spheres \(S_I\) are endowed with their standard spin structures;
\item the Lie group \(G\) admits a spin structure compatible with the chosen
metric \(g_G\).
\end{itemize}
Then every finite local model
\[
S_I\times G^I
\]
admits a natural local spin structure.
\end{proposition}

\begin{proof}
The tangent bundle of the local model splits as
\[
TS_I
\oplus
\bigoplus_{i\in I}TG.
\]

By assumption, both \(TS_I\) and \(TG\) admit spin structures. Since finite
direct sums of spin vector bundles again admit spin structures, the total
horizontal geometry carries a local spin structure.
\end{proof}

\subsubsection{Compatibility under gluing}

\begin{proposition}
Assume that the transition maps between finite local models preserve the local
spin structures. Then the local spin data glue into a global diffeological
spin structure on the horizontal pseudo-bundle.
\end{proposition}

\begin{proof}
The transition maps preserve the horizontal metric and identify the local
orthonormal geometries on overlaps. By assumption, the corresponding spin
lifts are compatible under these identifications. Hence the local spin
structures glue.
\end{proof}

\begin{remark}
This compatibility is automatic only if the transition maps lift coherently to
the spin coverings. In general, this constitutes the global obstruction to the
existence of a spin structure.
\end{remark}

\subsubsection{Spinor pseudo-bundle}

\begin{definition}
Let
\[
P_{\mathrm{Spin}}(\mathcal H)
\]
be a chosen spin lift of the horizontal orthogonal pseudo-bundle. The
associated spinor pseudo-bundle is
\[
\mathcal S
=
P_{\mathrm{Spin}}(\mathcal H)\times_\rho \Sigma,
\]
where
\[
\rho:\operatorname{Spin}(r)\to \operatorname{End}(\Sigma)
\]
is a spin representation.
\end{definition}

\begin{proposition}
The spinor pseudo-bundle \(\mathcal S\) carries a natural Clifford action
\[
\mathrm{Cl}(\mathcal H)\curvearrowright \mathcal S.
\]
\end{proposition}

\begin{proof}
This is the standard Clifford action induced by the spin representation on
each local model. Compatibility under gluing follows from compatibility of
the spin transition maps.
\end{proof}

\subsubsection{Spin connections}

\begin{proposition}
Assume that the horizontal metric admits a compatible metric connection
\[
\nabla^{\mathcal H}.
\]
Then the spinor pseudo-bundle \(\mathcal S\) carries an induced
Clifford-compatible spin connection
\[
\nabla^{\mathcal S}.
\]
\end{proposition}

\begin{proof}
Locally, the metric connection on the orthogonal bundle lifts to the spin
bundle through the differential of the covering
\[
\operatorname{Spin}(r)\to SO(r).
\]
The associated spin representation then induces a connection on the spinor
bundle. Compatibility with Clifford multiplication follows from the standard
spinorial construction.
\end{proof}

\subsubsection{Spin Dirac operator}

We now define the spin Dirac operator.

\begin{definition}
Assume first that the local horizontal geometry has finite rank. Let
\[
(e_a)
\]
be a local orthonormal frame of the horizontal geometry. The local spin Dirac
operator is
\[
D_{\mathrm{spin}}
=
\sum_a c(e_a)\nabla^{\mathcal S}_{e_a}.
\]
\end{definition}

\begin{proposition}
The local operator \(D_{\mathrm{spin}}\) is independent of the choice of local
orthonormal frame.
\end{proposition}

\begin{proof}
The proof is identical to the frame-independence proof for the general
Dirac-type operator. Orthogonal changes of frame preserve the Clifford
contraction because the orthogonal matrix coefficients satisfy
\[
\sum_a A_{ab}A_{ac}
=
\delta_{bc}.
\]
\end{proof}

\begin{proposition}
Assume that the local spin structures, spin connections, and Clifford actions
glue compatibly. Then the local spin Dirac operators define a global
diffeological spin Dirac operator.
\end{proposition}

\begin{proof}
The local expressions are frame-independent and compatible on overlaps by
construction. Hence they glue to a global operator.
\end{proof}

\subsubsection{Infinite-dimensional case}

\begin{remark}
If \(G\) is infinite-dimensional, the existence of spin structures becomes
substantially more delicate. In particular:
\begin{itemize}
\item infinite-dimensional Clifford algebras require Hilbert or convenient
completions;
\item spin representations may exist only for restricted orthogonal groups;
\item the corresponding Dirac operator requires additional convergence and
domain assumptions.
\end{itemize}
For this reason, the preceding spinorial construction should primarily be
understood on finite-rank local models unless additional analytic data are
specified.
\end{remark}

\subsubsection{Degeneracy directions}

\begin{remark}
The spin structure depends only on the effective horizontal pseudo-bundle
\[
\mathcal H
=
TE^{\mathrm{sph}}(G)/\mathcal K,
\]
and not on the degenerate directions contained in \(\mathcal K\).
\end{remark}

\begin{remark}
Consequently, the spin Dirac operator acts on the effective horizontal
geometry rather than on the full tangent pseudo-bundle of
\[
E^{\mathrm{sph}}(G).
\]
\end{remark}

\subsection{Summary}

Under the finite-rank or admissible analytic assumptions discussed above, we
have constructed:
\begin{itemize}
\item a Clifford structure on the horizontal pseudo-bundle;
\item a compatible metric connection;
\item local and global Dirac-type operators;
\item local spin structures and spinor pseudo-bundles;
\item a spin Dirac operator compatible with the horizontal geometry.
\end{itemize}

These constructions provide the differential-geometric framework required for
the Hodge- and Dirac-type theories developed in the sequel.
\section{Twisted structures and gerbes}

\subsection{$\mathbb{Z}_2$-local systems}

Let $X = E^{\mathrm{sph}}(G)$. Let
\[
\alpha \in H^1(X;\mathbb{Z}_2)
\]
be a $\mathbb{Z}_2$-valued cocycle.

\begin{definition}
The cocycle $\alpha$ defines a real line bundle $L_\alpha$ (or, more generally, a local system), with transition functions taking values in $\{\pm 1\}$.
\end{definition}

\begin{remark}
This twisting encodes sign changes compatible with the projective structure introduced in Section~3.
\end{remark}

\subsection{Twisted differential forms}

\begin{definition}
A twisted differential form is a differential form with values in $L_\alpha$:
\[
\Omega^k(X;L_\alpha).
\]
\end{definition}

\begin{proposition}
The exterior differential extends to twisted forms:
\[
d : \Omega^k(X;L_\alpha) \to \Omega^{k+1}(X;L_\alpha).
\]
\end{proposition}

\subsection{Twisted Dirac operator}

\begin{definition}
The twisted Dirac operator
\[
D_\alpha : \Omega^\bullet(X;L_\alpha) \to \Omega^\bullet(X;L_\alpha)
\]
is defined by the same local expression as $D$, acting on sections of the twisted bundle.
\end{definition}

\begin{remark}
Locally, the twist is invisible, and $D_\alpha$ coincides with $D$. The twisting appears in the transition functions.
\end{remark}

\subsection{Relation with extensions of Lie groups}

The appearance of $\mathbb{Z}_2$-twists is closely related to extensions of Lie groups.

\begin{remark}
In the framework developed by Neeb~\cite{Neeb2007}, such twists correspond to nontrivial extensions and cocycles in group cohomology.
\end{remark}

\subsection{Gerbe interpretation}

The twisted structures considered above can be interpreted in terms of higher geometric objects.

\begin{definition}
A gerbe on $X$ is, roughly speaking, a sheaf of groupoids with a band given by an abelian group (typically $\mathbb{U}(1)$ or $\mathbb{Z}_2$).
\end{definition}

\begin{remark}
The twisting by $\alpha$ can be viewed as the simplest instance of a gerbe, corresponding to a class in $H^2(X;\mathbb{Z}_2)$.
\end{remark}

\subsection{Higher curvature}

\begin{remark}
In the presence of connections on gerbes, one obtains curvature forms of degree $3$, generalizing the curvature of line bundles. Such structures have been studied in~\cite{BreenMessing2005}.
\end{remark}

\subsection{Compatibility with the geometric structures}

\begin{proposition}
The twisted structures are compatible with:
\begin{itemize}
\item the diffeological structure,
\item the Riemannian metric,
\item and the Dirac operator.
\end{itemize}
\end{proposition}

\subsection{Summary}

We have introduced:
\begin{itemize}
\item twisted differential forms,
\item a twisted Dirac operator,
\item and a connection with non-abelian extensions and gerbes.
\end{itemize}

\begin{remark}
These structures provide a higher-geometric extension of the framework developed in this article.
\end{remark}
\section{Classical structure groups}
\label{sec:classical-groups}

We describe the spherical Milnor construction for the classical groups
\[
O(n),\quad SO(n),\quad SU(n),\quad SO(1,n).
\]
These examples illustrate how the general diffeological construction specializes to familiar geometric structures.

\subsection{The orthogonal group \(O(n)\)}

Let
\[
G=O(n).
\]
The spherical Milnor space is
\[
E^{\mathrm{sph}}(O(n))
=
\left\{
(x_i,A_i)_{i\in\mathbb N}
\;\middle|\;
\sum_i x_i^2=1,\ A_i\in O(n),\ \text{finite support}
\right\}/\sim,
\]
where
\[
(x_i,A_i)\sim (x_i,A_i')
\quad\Longleftrightarrow\quad
A_i=A_i'
\text{ whenever }x_i\neq0.
\]

The group \(O(n)\) acts by left multiplication:
\[
B\cdot (x_i,A_i)=(x_i,BA_i).
\]

Since \(O(n)\) is compact, it carries a bi-invariant metric induced by
\[
\langle X,Y\rangle_{\mathfrak o(n)}
=
-\operatorname{tr}(XY),
\qquad
X,Y\in\mathfrak o(n).
\]
The barycentric metric on \(E^{\mathrm{sph}}(O(n))\) is locally
\[
g
=
\sum_i dx_i^2
+
\sum_i x_i^2\,\langle \theta_i,\theta_i\rangle_{\mathfrak o(n)},
\]
where \(\theta_i=A_i^{-1}dA_i\) is the Maurer--Cartan form on the \(i\)-th copy of \(O(n)\).

\begin{proposition}
The metric \(g\) is \(O(n)\)-invariant.
\end{proposition}

\begin{proof}
The coordinates \(x_i\) are unchanged by the left action. The Maurer--Cartan form is left-invariant and the inner product on \(\mathfrak o(n)\) is \(\operatorname{Ad}\)-invariant. Hence each term in the metric is invariant.
\end{proof}

Thus the construction gives a universal metric-type structure over
\[
B^{\mathrm{sph}}(O(n))
=
E^{\mathrm{sph}}(O(n))/O(n).
\]

\begin{remark}
Pulling this structure back by a classifying map
\[
f:X\to B^{\mathrm{sph}}(O(n))
\]
produces the corresponding geometric data on an \(O(n)\)-bundle over \(X\). This is the natural spherical analogue of the universal orthogonal frame bundle.
\end{remark}

\subsection{The special orthogonal group \(SO(n)\)}

For
\[
G=SO(n),
\]
one obtains
\[
E^{\mathrm{sph}}(SO(n))
\to
B^{\mathrm{sph}}(SO(n)).
\]
The Lie algebra is
\[
\mathfrak{so}(n)=\{X\in M_n(\mathbb R)\mid X^T+X=0\},
\]
with invariant inner product
\[
\langle X,Y\rangle_{\mathfrak{so}(n)}
=
-\operatorname{tr}(XY).
\]

The barycentric metric is
\[
g
=
\sum_i dx_i^2
+
\sum_i x_i^2\,\langle \theta_i,\theta_i\rangle_{\mathfrak{so}(n)}.
\]

\begin{proposition}
The spherical classifying space \(B^{\mathrm{sph}}(SO(n))\) carries the descended structures associated with oriented orthonormal bundles.
\end{proposition}

\begin{proof}
The proof is the same as for \(O(n)\), using the fact that \(SO(n)\) is compact and that the metric on \(\mathfrak{so}(n)\) is \(\operatorname{Ad}\)-invariant. Since the group is connected for \(n\geq2\), the quotient corresponds to the oriented case.
\end{proof}

\begin{remark}
The passage from \(O(n)\) to \(SO(n)\) corresponds to imposing an orientation. Thus \(B^{\mathrm{sph}}(SO(n))\) is the spherical classifying model for oriented rank \(n\) real vector bundles.
\end{remark}

\subsection{Spin refinement}

If \(n\geq2\), the double covering
\[
\operatorname{Spin}(n)\longrightarrow SO(n)
\]
gives a further lift of the construction:
\[
E^{\mathrm{sph}}(\operatorname{Spin}(n))
\to
B^{\mathrm{sph}}(\operatorname{Spin}(n)).
\]

\begin{remark}
The obstruction to lifting an \(SO(n)\)-classified object to the spin spherical classifying space is the usual second Stiefel--Whitney obstruction, expressed in this setting through the pullback of the corresponding universal class.
\end{remark}

When such a lift exists, the Clifford module and Dirac structures constructed in the previous sections become globally defined on the pulled-back bundle.

\subsection{The special unitary group \(SU(n)\)}

Let
\[
G=SU(n).
\]
Then
\[
\mathfrak{su}(n)
=
\{X\in M_n(\mathbb C)\mid X^*+X=0,\ \operatorname{tr}X=0\}.
\]

We use the invariant inner product
\[
\langle X,Y\rangle_{\mathfrak{su}(n)}
=
-\operatorname{Re}\operatorname{tr}(XY).
\]

The spherical Milnor space is
\[
E^{\mathrm{sph}}(SU(n))
=
\left\{
(x_i,U_i)_{i\in\mathbb N}
\;\middle|\;
\sum_i x_i^2=1,\ U_i\in SU(n),\ \text{finite support}
\right\}/\sim.
\]

The barycentric metric is
\[
g
=
\sum_i dx_i^2
+
\sum_i x_i^2\,\langle U_i^{-1}dU_i,U_i^{-1}dU_i\rangle_{\mathfrak{su}(n)}.
\]

\begin{proposition}
The metric and the associated Hodge-type and Dirac-type structures are \(SU(n)\)-invariant.
\end{proposition}

\begin{proof}
The group \(SU(n)\) is compact, and the inner product on \(\mathfrak{su}(n)\) is \(\operatorname{Ad}\)-invariant. The spherical coordinates are fixed by the left action, and the Maurer--Cartan forms transform equivariantly. Therefore the metric is invariant. The induced Hodge-type and Clifford structures are invariant whenever the chosen connection is \(SU(n)\)-compatible.
\end{proof}

\begin{remark}
The quotient
\[
B^{\mathrm{sph}}(SU(n))
\]
is the spherical classifying model for rank \(n\) complex vector bundles with trivial determinant.
\end{remark}

\subsection{Chern--Weil forms in the \(SU(n)\)-case}

Let
\[
\Theta
=
\sum_i x_i^2\,\theta_i
\]
be the barycentric connection form, with
\[
\theta_i=U_i^{-1}dU_i.
\]
Its curvature is
\[
F_\Theta
=
d\Theta+\frac12[\Theta,\Theta].
\]

For every invariant polynomial
\[
P\in \operatorname{Sym}^k(\mathfrak{su}(n)^*)^{SU(n)},
\]
the form
\[
P(F_\Theta,\ldots,F_\Theta)
\]
is basic and descends to
\[
B^{\mathrm{sph}}(SU(n)).
\]

\begin{remark}
These descended forms represent the usual Chern--Weil characteristic classes in the spherical diffeological model.
\end{remark}

\subsection{The Lorentz group \(SO(1,n)\)}

We now consider the non-compact group
\[
G=SO(1,n),
\]
defined as the identity component or full group preserving the quadratic form
\[
q(x)=-x_0^2+x_1^2+\cdots+x_n^2.
\]

Its Lie algebra is
\[
\mathfrak{so}(1,n)
=
\{X\in M_{n+1}(\mathbb R)\mid X^T\eta+\eta X=0\},
\]
where
\[
\eta=\operatorname{diag}(-1,1,\ldots,1).
\]

The spherical Milnor space is
\[
E^{\mathrm{sph}}(SO(1,n))
=
\left\{
(x_i,A_i)_{i\in\mathbb N}
\;\middle|\;
\sum_i x_i^2=1,\ A_i\in SO(1,n),\ \text{finite support}
\right\}/\sim.
\]

The left action is again
\[
B\cdot (x_i,A_i)=(x_i,BA_i).
\]

Unlike the compact cases, \(SO(1,n)\) does not carry a positive definite bi-invariant Riemannian metric. The Killing form is non-degenerate but indefinite. Hence the natural invariant structure is pseudo-Riemannian rather than Riemannian.

Let
\[
\kappa(X,Y)=\operatorname{tr}(\operatorname{ad}_X\operatorname{ad}_Y)
\]
be the Killing form. The barycentric pseudo-metric is
\[
g
=
\sum_i dx_i^2
+
\sum_i x_i^2\,\kappa(\theta_i,\theta_i).
\]

\begin{proposition}
The barycentric pseudo-metric on \(E^{\mathrm{sph}}(SO(1,n))\) is \(SO(1,n)\)-invariant.
\end{proposition}

\begin{proof}
The Killing form is \(\operatorname{Ad}\)-invariant, and the Maurer--Cartan forms transform by the adjoint action. The spherical part is unchanged by the group action. Therefore the pseudo-metric is invariant.
\end{proof}

\begin{remark}
Because the metric is indefinite, the associated Hodge and Dirac theories must be interpreted in a pseudo-Riemannian or hyperbolic sense. This distinguishes the Lorentzian case from the compact cases \(O(n)\), \(SO(n)\), and \(SU(n)\).
\end{remark}

\subsection{Cartan decomposition and positive metrics for \(SO(1,n)\)}

If a positive definite metric is desired, one may use a Cartan decomposition
\[
\mathfrak{so}(1,n)=\mathfrak k\oplus\mathfrak p,
\]
where
\[
\mathfrak k\simeq \mathfrak{so}(n)
\]
is the Lie algebra of a maximal compact subgroup \(SO(n)\).

The Cartan involution \(\theta\) defines a positive definite inner product
\[
\langle X,Y\rangle_\theta
=
-\kappa(X,\theta Y).
\]

This gives a left-invariant Riemannian metric on \(SO(1,n)\), hence a barycentric Riemannian metric on
\[
E^{\mathrm{sph}}(SO(1,n)).
\]

\begin{proposition}
The Cartan metric is left-invariant but not bi-invariant in general.
\end{proposition}

\begin{proof}
It is constructed by translating an inner product on the Lie algebra by left multiplication, hence it is left-invariant. It is not generally \(\operatorname{Ad}\)-invariant because \(SO(1,n)\) is non-compact and semisimple, and admits no positive definite bi-invariant metric.
\end{proof}

\begin{remark}
Thus, in the Lorentzian case, there are two distinct choices:
\begin{itemize}
\item an \(SO(1,n)\)-invariant pseudo-Riemannian structure from the Killing form;
\item a positive definite left-invariant Riemannian structure from the Cartan decomposition, which generally has weaker equivariance properties.
\end{itemize}
\end{remark}

\subsection{Lorentzian spin structures}

The double cover
\[
\operatorname{Spin}(1,n)\to SO(1,n)
\]
gives a corresponding spherical model
\[
E^{\mathrm{sph}}(\operatorname{Spin}(1,n))
\to
B^{\mathrm{sph}}(\operatorname{Spin}(1,n)).
\]

\begin{remark}
A lift from \(SO(1,n)\) to \(\operatorname{Spin}(1,n)\) gives the appropriate setting for Lorentzian spinor modules and Lorentzian Dirac operators. The resulting Dirac theory is hyperbolic rather than elliptic.
\end{remark}

\subsection{Summary}

The compact groups
\[
O(n),\quad SO(n),\quad SU(n)
\]
fit directly into the spherical Milnor framework with invariant positive definite metrics. Their spherical classifying spaces carry descended Riemannian, Hodge-type, and Dirac-type structures.

The Lorentz group
\[
SO(1,n)
\]
requires a distinction between invariant pseudo-Riemannian structures and positive definite Cartan metrics. This example shows that the spherical Milnor construction is flexible enough to include both compact gauge-type groups and non-compact groups arising in pseudo-Riemannian geometry.
\section{Classifying spaces of diffeomorphism groups and smooth fiber bundles}
\label{sec:diff-classifying}

Let \(M\) be a compact smooth manifold. We denote by
\[
Diff(M)
\]
the diffeomorphism group of \(M\), endowed either with its Fréchet Lie-group
structure when available, or with its natural diffeological structure.

The purpose of this section is to relate the spherical Milnor constructions of
the present article with the geometry of smooth fiber bundles and with the
descent problem for geometric operators.

\subsection{Classifying spaces of diffeomorphism groups}

The classifying space
\[
BDiff(M)
\]
classifies smooth fiber bundles with fiber \(M\). More precisely, if \(X\) is
a sufficiently regular smooth or diffeological space, then isomorphism classes
of smooth \(M\)-bundles
\[
M\longrightarrow E\longrightarrow X
\]
are represented by homotopy classes of classifying maps
\[
f:X\longrightarrow BDiff(M).
\]

The associated bundle is obtained by pullback from the universal bundle
\[
EDiff(M)\longrightarrow BDiff(M),
\]
namely
\[
E_f=f^*EDiff(M).
\]

\begin{remark}
In the diffeological setting, one should interpret the preceding statement in
the framework of diffeological principal bundles and diffeological homotopy
classes.
\end{remark}

\subsection{Spherical Milnor model}

Replacing the standard Milnor construction by the spherical construction of the
present article yields the universal pseudo-bundle
\[
E^{\mathrm{sph}}(Diff(M))
\longrightarrow
B^{\mathrm{sph}}(Diff(M)).
\]

A classifying map
\[
f:X\longrightarrow B^{\mathrm{sph}}(Diff(M))
\]
then determines a smooth \(M\)-bundle
\[
M\longrightarrow \mathcal E_f\longrightarrow X.
\]

\begin{remark}
The spherical construction does not modify the classification principle itself.
Rather, it enriches the universal geometry by additional structures:
\begin{itemize}
\item barycentric metrics;
\item horizontal pseudo-bundles;
\item Clifford structures;
\item Dirac-type operators;
\item and higher geometric structures associated with the spherical model.
\end{itemize}
\end{remark}

\subsection{Vertical geometry of the associated bundle}

Let
\[
\pi:\mathcal E_f\longrightarrow X
\]
be the bundle associated with a classifying map.

The vertical tangent pseudo-bundle is
\[
T^{\mathrm{vert}}\mathcal E_f
:=
\ker(d\pi).
\]

Locally, the bundle is obtained from transition functions valued in
\[
Diff(M).
\]
Therefore any geometric structure on the fiber \(M\) descends fiberwise only
if it is preserved by these transition functions.

Typical examples include:
\begin{itemize}
\item orientations;
\item Riemannian metrics;
\item symplectic structures;
\item spin structures;
\item Clifford modules;
\item Dirac operators.
\end{itemize}

\subsection{Reduction of the structure group}

Let \(\mathcal S\) be a geometric structure on \(M\). We denote by
\[
Diff(M,\mathcal S)
\subset
Diff(M)
\]
the subgroup preserving \(\mathcal S\).

Typical examples are:
\[
Diff^+(M),
\qquad
Isom(M,g),
\qquad
Diff^{\mathrm{spin}}(M).
\]

\begin{proposition}
A smooth \(M\)-bundle
\[
M\longrightarrow \mathcal E\longrightarrow X
\]
admits a fiberwise structure of type \(\mathcal S\) if and only if its
classifying map
\[
f:X\longrightarrow BDiff(M)
\]
lifts, up to homotopy, to
\[
BDiff(M,\mathcal S).
\]
\end{proposition}

\begin{proof}
The bundle is locally glued by transition functions valued in
\[
Diff(M).
\]
A fiberwise structure of type \(\mathcal S\) exists precisely when the
transition functions preserve \(\mathcal S\), that is, when they take values
in the subgroup
\[
Diff(M,\mathcal S).
\]

This is exactly the datum of a reduction of the structure group from
\[
Diff(M)
\]
to
\[
Diff(M,\mathcal S).
\]

Such reductions are classified by lifts of the classifying map to the
classifying space of the subgroup.
\end{proof}

\subsection{Fiberwise spin structures and Dirac operators}

Assume now that:
\begin{itemize}
\item \(M\) is spin;
\item the structure group reduces to
\[
Diff^{\mathrm{spin}}(M).
\]
\end{itemize}

Then the vertical tangent bundle
\[
T^{\mathrm{vert}}\mathcal E_f
\]
inherits a fiberwise spin structure.

\begin{proposition}
Under the preceding assumptions, there exists a vertical spinor pseudo-bundle
\[
S^{\mathrm{vert}}\mathcal E_f
\]
together with a family of vertical Dirac operators
\[
D_x:
\Gamma(S_{\mathcal E_x})
\longrightarrow
\Gamma(S_{\mathcal E_x}),
\qquad x\in X.
\]
\end{proposition}

\begin{proof}
The reduction of the structure group to
\[
Diff^{\mathrm{spin}}(M)
\]
ensures that the spin structures on the fibers are preserved by the transition
functions. Hence the local spinor bundles glue into a global vertical spinor
pseudo-bundle.

The local vertical Dirac operators are compatible under these spinorial
transition maps and therefore define a smooth family.
\end{proof}

\begin{remark}
The preceding construction is entirely fiberwise. It does not require a global
spin structure on the total space \(\mathcal E_f\).
\end{remark}

\subsection{Non-invariant geometric structures}

We now turn to the case where the geometric structure is not preserved by the
full diffeomorphism group.

Let
\[
D_M
\]
be a Dirac-type operator on \(M\). For
\[
\varphi\in Diff(M),
\]
define the transformed operator
\[
\varphi_*D_M\varphi_*^{-1}.
\]

\begin{definition}
The defect of invariance of the operator \(D_M\) under the diffeomorphism
\(\varphi\) is
\[
\mathcal D_D(\varphi)
:=
\varphi_*D_M\varphi_*^{-1}-D_M.
\]
\end{definition}

\begin{proposition}
The defect vanishes if and only if \(\varphi\) preserves the Dirac structure.
\end{proposition}

\begin{proof}
The condition
\[
\mathcal D_D(\varphi)=0
\]
is exactly
\[
\varphi_*D_M\varphi_*^{-1}=D_M,
\]
which means that the operator is invariant under the action of
\(\varphi\).
\end{proof}

\subsection{Infinitesimal defect}

Let
\[
X\in\mathfrak X(M)
\]
be a smooth vector field generating a local flow
\[
(\varphi_t).
\]

\begin{definition}
The infinitesimal defect of \(D_M\) along \(X\) is
\[
\mathcal D_D(X)
:=
[\mathcal L_X,D_M],
\]
where \(\mathcal L_X\) denotes the Lie derivative.
\end{definition}

\begin{proposition}
The infinitesimal defect vanishes if and only if the flow generated by \(X\)
preserves the Dirac structure infinitesimally.
\end{proposition}

\begin{proof}
Differentiating the identity
\[
\varphi_t^*D_M(\varphi_t^{-1})^*
=
D_M
\]
at \(t=0\) yields
\[
[\mathcal L_X,D_M]=0.
\]
Conversely, vanishing of the commutator implies infinitesimal invariance along
the flow generated by \(X\).
\end{proof}

\subsection{Defect bundles and characteristic classes}

The preceding constructions suggest that the failure of invariance of a
geometric structure should itself define a geometric object over the
classifying space.

Roughly speaking, the defect measures the obstruction to descending a chosen
geometric operator from the universal fiber \(M\) to the associated bundle.

\begin{definition}
A defect bundle associated with a geometric operator \(D_M\) is a bundle-like
object encoding the variation of the transformed operators
\[
\varphi_*D_M\varphi_*^{-1}
\]
along the structure group.
\end{definition}

\begin{remark}
The precise analytic realization of such defect bundles depends strongly on
the operator class under consideration:
\begin{itemize}
\item finite-dimensional elliptic operators;
\item pseudodifferential operators;
\item infinite-dimensional Dirac operators;
\item current-group constructions;
\item determinant and gerbe-type extensions.
\end{itemize}
\end{remark}

\begin{remark}
In favorable situations, the defect construction gives rise to characteristic
classes measuring the obstruction to reducing the structure group to the
subgroup preserving the chosen geometric structure.
\end{remark}

\subsection{Outlook}

The spherical Milnor model therefore provides a natural framework for:
\begin{itemize}
\item universal geometric structures on classifying spaces;
\item fiberwise Clifford and Dirac geometries;
\item obstruction and defect theories;
\item determinant-type constructions;
\item and higher geometric structures associated with diffeomorphism groups.
\end{itemize}

These constructions will be further developed in subsequent works devoted to:
\begin{itemize}
\item current groups and restricted symmetry groups;
\item determinant bundles and Schwinger cocycles;
\item gerbe-theoretic formulations;
\item and applications to integrable systems and geometric field theories.
\end{itemize}


\appendix

\section{Relation with Fisher--Rao type structures}

\subsection{Fisher--Rao metric on the simplex}

Let
\[
\Delta_I
=
\left\{ (t_i)_{i\in I} \in \mathbb{R}^I \;\middle|\; t_i \ge 0,\ \sum_{i\in I} t_i = 1 \right\}.
\]

The Fisher--Rao metric on $\Delta_I$ is given by
\[
g_{\mathrm{FR}}
=
\sum_{i\in I} \frac{dt_i^2}{t_i},
\]
defined on the interior of the simplex.

\begin{remark}
This metric is singular at the boundary $t_i=0$.
\end{remark}

\subsection{Spherical parametrization}

Consider the change of variables
\[
t_i = x_i^2,
\qquad
\sum_i x_i^2 = 1.
\]

\begin{proposition}
Under this transformation, the Fisher--Rao metric becomes
\[
g_{\mathrm{FR}} = 4 \sum_i dx_i^2.
\]
\end{proposition}

\begin{proof}
We compute
\[
dt_i = 2x_i dx_i,
\]
hence
\[
\frac{dt_i^2}{t_i}
=
\frac{4x_i^2 dx_i^2}{x_i^2}
=
4 dx_i^2.
\]
Summing over $i$ gives the result.
\end{proof}

\begin{remark}
Thus the Fisher--Rao metric becomes a multiple of the Euclidean metric on the sphere.
\end{remark}

\subsection{Interpretation}

\begin{remark}
This shows that the spherical normalization removes the singularity of the Fisher--Rao metric at the boundary of the simplex.
\end{remark}

\begin{remark}
The degeneracy observed in the spherical Milnor space is therefore not due to the metric itself, but to the quotient structure on the group coordinates.
\end{remark}

\subsection{Extension to the Milnor setting}

In the spherical Milnor space, the metric takes the form
\[
g = \sum_i dx_i^2 + \sum_i x_i^2 \, g_G.
\]

\begin{remark}
The first term corresponds to the Fisher--Rao metric in spherical coordinates, while the second term introduces a coupling with the group geometry.
\end{remark}

\subsection{Second-order effects}

\begin{proposition}
Let $x_i(t) = t$ near $t=0$. Then
\[
t_i(t) = x_i(t)^2 = t^2,
\]
and
\[
\frac{dt_i}{dt}(0) = 0.
\]
\end{proposition}

\begin{remark}
This shows that boundary directions are invisible at first order in Fisher--Rao coordinates.
\end{remark}

\begin{remark}
This phenomenon matches the second-order behavior described in Section~5.
\end{remark}

\subsection{Conclusion}

\begin{remark}
The spherical Milnor construction can be interpreted as a geometric desingularization of Fisher--Rao type metrics, combined with a coupling to group-valued degrees of freedom.
\end{remark}
\subsection{Outlook}

\begin{remark}
It would be interesting to further explore the relation between information geometry and the geometric structures developed in this article, in particular in connection with statistical or probabilistic interpretations.
\end{remark}

\vskip 12pt

\paragraph{Conflict of interest statement:} The author declares no conflict of interest.

\vskip 12pt

\paragraph{\bf Acknowledgements:} J.-P.M  thanks the France 2030 framework programme Centre Henri Lebesgue ANR-11-LABX-0020-01 
for creating an attractive mathematical environment.

\vskip 12pt

\paragraph{\bf Declaration of generative AI and AI-assisted technologies in the writing process}

During the preparation of this work the author used ChatGPT and Mistral AI in order to smoothen the expression in English. After using this tool/service, the author reviewed and edited the content as needed and takes full responsibility for the content of the publication.
\begin{thebibliography}{99}
%	\bibitem{AFHKL1986}	S. Albeverio, J. E. Fenstad, R. Hoegh-Krohn, T. Lindstrom,  \emph{Nonstandard Methods in Stochastic Analysis and Mathematical Physics} Dover (1986). 
    
\bibitem{AmariNagaoka2000}
S.-I.~Amari and H.~Nagaoka,
\newblock {\em Methods of Information Geometry},
\newblock Translations of Mathematical Monographs, vol.~191,
\newblock American Mathematical Society, 2000.

\bibitem{AyJostLeSchwachhofer2017}
N.~Ay, J.~Jost, H.~V.~L\^e, and L.~Schwachh\"ofer,
\newblock {\em Information Geometry},
\newblock Springer, 2017.

	\bibitem{BIgKWa2014} Batubenge, A.; Karshon, Y.; Watts, J.A.; Diffeological, Fr\"olicher, and differential spaces. \emph{Rocky Mountain J. Math.} 55 (2025), No. 3, 637-670.

\bibitem{BaezHoffnungRogers2010}
J.~C.~Baez, A.~E.~Hoffnung,
\newblock Convenient categories of smooth spaces,
\newblock {\em Trans. Amer. Math. Soc.} 363 (2011), 5789--5825.
	
\bibitem{BaezLauda2004}
J.~C.~Baez and A.~D.~Lauda,
\newblock Higher-dimensional algebra V: 2-groups,
\newblock {\em Theory Appl. Categ.} 12 (2004), 423--491.

\bibitem{BottTu1982}
R.~Bott and L.~W.~Tu,
\newblock {\em Differential Forms in Algebraic Topology},
\newblock Springer, 1982.

\bibitem{Breen2008}
L.~Breen,
\newblock Differential geometry of gerbes and differential forms,
\newblock arXiv:0802.1833.

\bibitem{BreenMessing2005}
L.~Breen and W.~Messing,
\newblock Differential geometry of gerbes,
\newblock {\em Adv. Math.} 198 (2005), 732--846.

\bibitem{CareyMickelssonMurray1997}
A.~L.~Carey, J.~Mickelsson, and M.~K.~Murray,
\newblock Index theory, gerbes, and Hamiltonian quantization,
\newblock {\em Comm. Math. Phys.} 183 (1997), 707--722.

\bibitem{ChristensenSinnamonWu2014}
J.~D.~Christensen, G.~Sinnamon, and E.~Wu,
\newblock The $D$-topology for diffeological spaces,
\newblock {\em Pacific J. Math.} 272 (2014), 87--110.

\bibitem{ChristensenWu2016}
J.~D.~Christensen and E.~Wu,
\newblock Tangent spaces and tangent bundles for diffeological spaces,
\newblock {\em Cah. Topol. G\'eom. Diff\'erentielle Cat\'eg.} 57 (2016), 3--50.

\bibitem{ChristensenWuSmooth}
J.~D.~Christensen and E.~Wu,
\newblock Smooth classifying spaces,
\newblock {\em Cah. Topol. G\'eom. Diff\'erentielle Cat\'eg.},
\newblock (à préciser volume/année si vous voulez être ultra précis).

	\bibitem{DN2007-1}
		D. Dugmore, P. Ntumba, On tangent cones of Fr\"olicher spaces.
		\textit{Quaetiones mathematicae} \textbf{30} (2007) no. 1, 67--83.
	%	\bibitem{DN2007-2} D. Dugmore, P. Ntumba, Cofibrations in the category of Fr\"olicher spaces: part I. \textit{Homotopy, homology and applications} \textbf{9} (2007) no. 2, 413--444.

\bibitem{Freed1995}
D.~S.~Freed,
\newblock Classical Chern--Simons theory. Part 1,
\newblock {\em Adv. Math.} 113 (1995), 237--303.

\bibitem{GMW2023} N. Goldammer, J.-P. Magnot, K. Welker, On diffeologies from infinite dimensional geometry to PDE constrained optimization. \emph{Contemp. Math.} {\bf 794} 1-48 (2024). 

\bibitem{GW2020} N. Goldammer, K. Welker,  Towards optimization techniques on diffeological spaces by generalizing Riemannian concepts. \emph{Appl. Math. Optim.} {\bf 93} (2026) No. 1, Paper No. 5, 31 p. 

\bibitem{He1995} G. Hector,
		G\'eom\'etrie et topologie des espaces diff\'eologiques.
		, in: \textit{Analysis and Geometry in Foliated Manifolds
			(Santiago de Compostela, 1994)}, World Sci. Publishing,  55--80 (1995).

\bibitem{IglesiasZemmour2013}
P.~Iglesias-Zemmour,
\newblock {\em Diffeology},
\newblock American Mathematical Society, 2013.

	%\bibitem{KW2016} Y. Karshon, J. Watts,  Basic forms and orbit spaces: a diffeological approach. \emph{SIGMA, Symmetry Integrability Geom. Methods Appl.} {\bf 12} (2006), Paper 026, 19 p.
	%\bibitem{KW} B. Khesin, R. Wendt, \textit{The Geometry of Infinite-Dimensional groups}. 
	%Springer-Verlag (2009).
		% \bibitem{KLMP2013} B. Khesin,  J. Lenells, G. Misiolek, S. C. Preston,  Geometry of diffeomorphism groups, complete integrability and optimal transport. \emph{Pure and Applied
%Math. Quarterly} {\bf 9} (2013) no. 2, 292--332.
		\bibitem{KM} A. Kriegl, P. W. Michor, \textit{The convenient setting
		for global analysis.} Math. surveys and monographs \textbf{53}, American
	Mathematical society, Providence, USA. (2000).

    	\bibitem{Les} J. Leslie,  On a Diffeological Group Realization of
	certain Generalized symmetrizable Kac-Moody Lie Algebras. \textit{J.
		Lie Theory} \textbf{13} (2003), 427--442.
    
    \bibitem{Ma2006-3} J.-P. Magnot,  Diff\'eologie sur le fibr\'e d'holonomie d'une connexion en dimension infinie.
	\textit{C. R. Math. Acad. Sci., Soc. R. Can.} \textbf{28} (2006) no. 4, 121--128.

    	\bibitem{Ma2013} J.-P. Magnot,  Ambrose-Singer theorem on diffeological bundles and complete integrability of the KP equation. \textit{Int. J. Geom. Meth. Mod. Phys.} \textbf{ 10} (2013) no. 9, Article ID 1350043, 31 p.  
\bibitem{MagnotWatts2017}
J.-P.~Magnot and J.~Watts,
\newblock The diffeology of Milnor's classifying space,
\newblock {\em Topology Appl.} 232 (2017), 189--213.

\bibitem{Magnotloops2019} J. P. Magot,
\newblock Remarks on the geometry and the topology of the loop spaces $H^s(S^1,N),$
for $s \leq 1/2.$
\newblock \emph{International Journal of Maps in Mathematics} {\bf 2} (2019) no1, pp.14 - 37.

\bibitem{Magnot-Cauchy} J.-P.~Magnot,
\newblock  On the differential geometry of numerical schemes and weak solutions of functional equations. 
\newblock \emph{Nonlinearity} {\bf 33} (2020), No. 12, 6835-6867. 
\bibitem{MagnotConnections}
J.-P.~Magnot,
\newblock  On the geometry of diffeological vector pseudobundles and infinite dimensional vector bundles: automorphisms, connections and covariant derivatives. 
\newblock \emph{Carpathian J. Math.} {\bf 41} (2025), No. 3, 687-710. 



\bibitem{Milnor1956I}
J.~Milnor,
\newblock Construction of universal bundles. I,
\newblock {\em Ann. of Math.} 63 (1956), 272--284.

\bibitem{Milnor1956II}
J.~Milnor,
\newblock Construction of universal bundles. II,
\newblock {\em Ann. of Math.} 63 (1956), 430--436.

\bibitem{Murray1996}
M.~K.~Murray,
\newblock Bundle gerbes,
\newblock {\em J. London Math. Soc.} 54 (1996), 403--416.

\bibitem{Neeb2004}
K.-H.~Neeb,
\newblock Abelian extensions of infinite-dimensional Lie groups,
\newblock {\em Travaux Math.} 15 (2004), 69--194.

\bibitem{Neeb2007}
K.-H.~Neeb,
\newblock Non-abelian extensions of infinite-dimensional Lie groups,
\newblock {\em Ann. Inst. Fourier} 57 (2007), 209--271.

\bibitem{NikolausWaldorf2013}
T.~Nikolaus and K.~Waldorf,
\newblock Four equivalent versions of non-abelian gerbes,
\newblock {\em Pacific J. Math.} 264 (2013), 355--420.

\bibitem{Pervova2016}
E. Pervova,
\newblock Diffeological vector pseudo-bundles. 
\newblock \emph{Topology Appl.} {\bf 202} (2016), 269-300.

\bibitem{Pervova2017}
E. Pervova,
\newblock Diffeological gluing of vector pseudo-bundles and pseudo-metrics on them. 
\newblock \emph{Topology Appl.} {\bf 220} (2017), 65-99. 

\bibitem{Pervova2019}
E. Pervova, 
\newblock Diffeological Clifford algebras and pseudo-bundles of Clifford modules. \newblock \emph{Linear Multilinear Algebra} {\bf 67} (2019), No. 9, 1785-1828. 

%\bibitem{PressleySegal1986}
%A.~Pressley and G.~Segal,
%\newblock {\em Loop Groups},
%\newblock Oxford Univ. Press, 1986.

\bibitem{Rao1945}
C.~R.~Rao,
\newblock Information and the accuracy attainable in the estimation of statistical parameters,
\newblock {\em Bull. Calcutta Math. Soc.} 37 (1945), 81--91.

\bibitem{Savelyev2026}
Y. Savelyev,
\newblock Smooth simplicial sets and universal Chern-Weil for infinite-dimensional groups. 
\newblock \emph{J. Lond. Math. Soc., II. Ser.} 113 no. 1 (2026), Article ID e70417, 65 p. 

\bibitem{Scalbi2024}
N.~T.~Scalbi,
\newblock Differential cohomology as diffeological homotopy theory,
\newblock arXiv:2408.02593.

\bibitem{Souriau1980}
J.-M.~Souriau,
\newblock Groupes diff\'erentiels,
\newblock Lecture Notes in Math., vol.~836, Springer, 1980.

	\bibitem{Sou} J.-M. Souriau, Un algorithme g\'en\'erateur de structures quantiques; 
		\textit{Ast\'erisque}, Hors S\'erie (1985), 341-399.

\bibitem{Stacey2011}
A.~Stacey,
\newblock Comparative smootheology,
\newblock {\em Theory Appl. Categ.} 25 (2011), 64--117.

\bibitem{Waldorf2007}
K.~Waldorf,
\newblock More morphisms between bundle gerbes,
\newblock {\em Theory Appl. Categ.} 18 (2007), 240--273.

\bibitem{Watts2012Thesis}
J.~Watts,
\newblock {\em Diffeologies, Differential Spaces, and Symplectic Geometry},
\newblock Ph.D. thesis, University of Toronto, 2012.

\bibitem{Wockel2008}
C.~Wockel,
\newblock Principal 2-bundles and their classifying spaces,
\newblock arXiv:0803.3692.

\end{thebibliography}
\end{document}
